% note.tex — 陽子どうしの衝突の NLO の QCD の補正：Drell--Yan、gg -> H（重いトップの有効理論）、MC@NLO（MS-bar、d = 4 - 2 eps）
% コンパイル: lualatex note.tex（目次と参照の解決のため2〜3回）
\documentclass[a4paper,10pt]{ltjsarticle}
\usepackage[margin=20mm]{geometry}
\usepackage[haranoaji,deluxe]{luatexja-preset}
\ltjsetparameter{jacharrange={-3}}
\usepackage{amsmath,amssymb,slashed}
\usepackage{booktabs,tabularx,array}
\usepackage[table]{xcolor}
\usepackage[font=small,labelfont=bf]{caption}
\usepackage{tikz}
\usetikzlibrary{arrows.meta,calc,positioning,decorations.markings,decorations.pathmorphing,shapes.geometric}
\usepackage{pgfplots}
\pgfplotsset{compat=1.18}
% the figures of step (3): categorical colours validated for colour-vision deficiency (dataviz palette slots 1-3), with line styles
\definecolor{mcblue}{HTML}{2A78D6}
\definecolor{ownorange}{HTML}{EB6834}
\definecolor{foaqua}{HTML}{1BAF7A}
\pgfplotsset{mcstyle/.style={width=7.2cm, tick label style={font=\scriptsize}, label style={font=\small}, title style={font=\small},
  legend style={font=\scriptsize, draw=none, fill=white, fill opacity=0.9, text opacity=1, cells={anchor=west}}, unbounded coords=jump}}
\usepackage{placeins}
\usepackage{xurl}
\usepackage{hyperref}
\hypersetup{colorlinks=true,linkcolor=blue!50!black,citecolor=green!40!black,urlcolor=blue!60!black,
  pdftitle={NLO QCD corrections in proton-proton collisions: Drell-Yan, gg -> H in the heavy-top EFT, MC@NLO},pdfauthor={}}

\setlength{\emergencystretch}{3em}
\renewcommand{\jsParagraphMark}{\textcolor{black!65}{\rule[-0.02\zw]{0.62\zw}{0.62\zw}}\hspace{0.35\zw}}
\newcommand{\dd}{\mathrm{d}}
\newcommand{\MSbar}{\overline{\mathrm{MS}}}
\newcommand{\MGaMC}{MG5\_aMC@NLO}
\newcommand{\re}{\mathrm{Re}\,}

\tikzset{
  photon/.style={thick,decorate,decoration={snake,amplitude=1.1pt,segment length=4.6pt}},
  gluon/.style={thick,decorate,decoration={coil,amplitude=2.3pt,segment length=3.6pt,aspect=0.75}},
  fermion/.style={thick,postaction={decorate},
    decoration={markings,mark=at position 0.58 with {\arrow{Stealth[length=1.9mm]}}}},
  vtx/.style={circle,fill=black,inner sep=1.3pt},
  blob/.style={ellipse,draw=black,thick,fill=black!12,minimum width=0.9cm,minimum height=1.25cm,inner sep=0pt},
  lab/.style={font=\small,align=center},
}

\title{陽子どうしの衝突の NLO の QCD の補正\\（Drell--Yan、重いトップの有効理論の $gg\to H$、MC@NLO、$\MSbar$、$d=4-2\epsilon$）}
\author{}
\date{2026年10月20日}

\begin{document}
\maketitle

\begin{abstract}
SM ノート~\cite{smnote}、CKM ノート~\cite{ckmnote}、NLO ノート~\cite{nlonote}、ヒッグスのノート~\cite{higgsnote} は、崩壊の NLO を実の放射まで含めて求めた。QCD ノート~\cite{qcdnote} は、$e^+e^-\to$ ハドロンと重いクォークの対生成の NLO を、QCD だけの模型で aMC@NLO と突き合わせた。本ノートでは、陽子の中のパートンが始状態にある色を持たない粒子の生成の NLO の QCD の補正を扱う（パートン、パートン分布、NLO の補正の考え方は、第~\ref{sec:primer} 節で学部の物理から説明する）。$\overline{\mathrm{MS}}$ の共線の対抗項でパートン分布に吸収する始状態の放射が、終状態だけの崩壊と違う点である。四つの手順の計画を立てる。(1) Drell--Yan（$pp\to\gamma^*/Z\to e^+e^-$、$W^\pm\to e\nu$）、(2) 重いトップの有効理論の $gg\to H$、(3) MC@NLO と Pythia 8、(4) 段階 6 の $Z'$ の NLO。手順 (1) では、$q\bar q\to V^*g$、$qg\to V^*q$ の実の放射と形状因子を $d$ 次元で $\epsilon$ について厳密に求め、$\MSbar$ の共線の対抗項を足すと極が記号として消え、$\ln\mu_R$ が落ち、係数関数が Altarelli--Ellis--Martinelli の $\MSbar$ の形（参照）と記号として一致することを確かめた。4次元の振幅の2乗は、本ノートの FeynArts の模型の $\gamma$、$Z$、$W$ の4過程と記号として一致した。係数関数を NNPDF4.0 NLO と自分で畳み込んだ断面積を、\MGaMC{} の固定次数の LO と NLO（\texttt{[QCD]}）と比べた。$Z$ の共鳴の付近、0.5~TeV より上、$W^\pm$ の全断面積と質量の窓、$m_{\ell\ell}$ の分布、ラピディティの分布（\ref{app:rap} の自分の $d$ 次元の式）、9通りの $(\mu_R,\mu_F)$ の重みの付け直しは、LO と NLO で統計の範囲で一致した（$Z$ の NLO の全断面積は $2064.54$~pb に対し $2064.78\pm0.22$~pb）。$K$ は $Z$、$W^+$、$W^-$ で 1.140、1.150、1.138（$q\bar q$ が $+29\%$、$qg$ が $-15\%$）、0.5~TeV より上で 1.188 で、スケールの範囲は LO の $\pm12\%$ から NLO の $^{+4}_{-7}\%$ に縮む。$W$ の NLO の全断面積は \MGaMC{} の積分ごとのばらつきが誤差の約2.4倍あり、その源は $m_{e\nu}<1$~GeV の正と負の重みの打ち消しだった。自分の積分では、Breit--Wigner の変数だけの求積が $W$ の裾を $10^{-3}$ 少なく見積もっていたことが、この比較で分かった。手順 (2) では、重いトップの有効理論の $gg\to H$ の $gg$、$qg$、$q\bar q$ の係数関数を $d$ 次元で求め、文献の形（参照）と記号として一致させた。NLOCT は次元5の演算子の頂点の R2 と紫外の対抗項を出せなかったので、まず HC\_NLO\_X0（FeynRules の模型のデータベース）を使った。その MadLoop の仮想の補正は、ウィルソン係数の補正を足すと自分の値と $2\times10^{-14}$ で一致したが、軽いクォークの紫外の対抗項の極が MG5 3.8.1 で落ちて aMC@NLO が止まるので、作業用の複製で直した。$\mu_R$ の重みの付け直しが $g_{Hgg}$ の中の $\alpha_s^2$ を動かさないこと、赤外で有限な $q\bar q\to Hg$ が \texttt{[QCD]} に入らないことも分かり、それらを考えに入れると LO と NLO の断面積は統計の範囲で一致した。$\mu=M_H$ で $K=2.29$ である。次に、$GGH$ の頂点の R2 と紫外の対抗項を FeynCalc で自分で導き（R2 は HC\_NLO\_X0 のもの（参照）と7つの係数がすべて一致した）、NLOCT の $H$ の頂点と入れ替えて自作の NLO の UFO を作った。MadLoop は直しなしで自分の値と $2\times10^{-14}$ で一致し、aMC@NLO の LO と NLO は、今度は $\mu_R$ を含む9通りの $(\mu_R,\mu_F)$ の重みの付け直しで、統計の範囲で一致した。手順 (3) では、$Z$ と $H$ を MC@NLO で作り、Pythia 8 でシャワーをかけた（パートンの段階）。全断面積は自分の NLO と統計の範囲で一致し（$-0.06\sigma$、$-0.45\sigma$）、ボソンのラピディティは NLO と数 \% で合った。大きい $p_T$ では、自分で4次元の実の放射から求めた LO の V$+$ジェットに近づき（150〜250~GeV で比 1.00、1.035）、小さい $p_T$ では固定次数の発散が有限な分布になった。その途中で、スケールの重みを切ると \MGaMC{} 3.8.1 のシャワーの段がメモリを使い尽くすことなど、ツールの問題を三つ見つけた。手順 (4) では、計算環境の記録の段階 6 の SSM の $Z'$ を NLO にした。手順 (1) の係数関数に $Z'$ の Born を入れるだけで $K$（ピークの付近で 1.18〜1.41）が求まり、QCD と $Z'$ だけの自作の NLO の UFO の MadLoop（$10^{-15}$）と aMC@NLO（統計の範囲）がこれを確かめた。段階 6 の事象を $K(m_{\ell\ell})$ で重み付けし直すと、SSM の質量の下限は LO の 5.04~TeV から 5.20~TeV（期待 5.16~TeV）に上がった（ATLAS は NNLO で 5.1~TeV、参照）。
\end{abstract}

\tableofcontents

\section{目的と問い}\label{sec:intro}

本ノートの問いは三つである。\textbf{第一に、始状態にパートンがある色を持たない粒子の生成の NLO の QCD の補正は、$d$ 次元の仮想の補正と実の放射に $\MSbar$ の共線の対抗項を足すと有限になり、$\ln\mu_R$ によらず、$\ln\mu_F$ の依存がパートン分布の DGLAP の進化と打ち消し合う形になるか。第二に、\MGaMC{} の FKS の引き算による固定次数の NLO（\texttt{[QCD]}）は、ハドロンの断面積でこれを再現するか。$\mu_R$、$\mu_F$ の重みの付け直しは正しいか。第三に、MC@NLO で Pythia 8 のシャワーをつなぐと、全断面積と分布はどう変わるか。}

第一の問いを立てる理由は、崩壊の NLO（NLO ノート、ヒッグスのノート）では終状態の赤外の発散が実と仮想の和で消えたのに対し、始状態の共線の発散は消えず、パートン分布の $\MSbar$ の繰り込みで吸収されるからである。その処方の規約（$d$ 次元のグルーオンの偏極の平均 $1/(2(1-\epsilon))$、$\overline{\mathrm{MS}}$ の極だけを引くこと）がパートン分布の決め方と合っていないと、ハドロンの断面積が有限部分でずれる。陽子の中のパートンの NLO は、LHC の信号と背景の規格化（$K$ 因子）の土台である。

第二の問いは、計算環境の記録~\cite{env} で \MGaMC{} の \texttt{[QCD]} を確かめた範囲に関わる。QCD ノートは、$e^+e^-\to u\bar u$、$t\bar t$ の $K$ 因子と、QCD だけの模型での $pp\to t\bar t$ を、自分の計算と突き合わせた。電弱の結合を持つ標準模型の過程（MadGraph 同梱の \texttt{loop\_sm}）、共鳴の Breit--Wigner、ハドロンの分布の重みの付け直しは、まだ確かめていない。段階 6 の $Z'$ の断面積は LO のままである。

第三の問いは、LHC の事象の生成に NLO を使うときの、MC@NLO の整合の確かめである。

求める量は四つである。(i) Drell--Yan の $\MSbar$ の係数関数 $\Delta_{q\bar q}$、$\Delta_{qg}$ とハドロンの $\dd\sigma/\dd Q^2$。(ii) 重いトップの有効理論の $gg\to H$ の係数関数とウィルソン係数の $O(\alpha_s)$ の補正。(iii) MC@NLO の全断面積と分布。(iv) $Z'$ の $K$ 因子。

本ノートの構成は次の通りである。第~\ref{sec:primer} 節で、パートンとその NLO の考え方を学部の物理から説明する（パートンになじみのない読者はここから読むとよい）。第~\ref{sec:setup} 節で模型と入力を定める。第~\ref{sec:results} 節で手計算の結果を述べる。第~\ref{sec:plan} 節でツールと突き合わせる計画を定める。第~\ref{sec:checks} 節に各手順の結果を、第~\ref{sec:discussion} 節にまとめを書く。導出の細部は\ref{app:dy} にある。

\section{陽子の中のパートン：学部の物理からの導入}\label{sec:primer}

この節は、量子力学、特殊相対論、電磁気学とクォーク模型の初歩（学部の物理）だけを前提に、本ノートで使う考え方と記号を説明する。第~\ref{sec:setup} 節から後の式は、すべてこの節の言葉で読める。最後に用語の対応表（表~\ref{tab:glossary}）を置く。

\subsection{陽子は何でできているか}\label{sec:proton}

陽子は、電荷 $+\frac23e$ の $u$ クォーク2個と $-\frac13e$ の $d$ クォーク1個（$uud$）でできている、とクォーク模型は教える。しかしこの3個の質量は合わせて約 9~MeV で、陽子の質量 938~MeV の $1\%$ にすぎない。残りは、クォークを結びつける強い相互作用の場のエネルギーである。強い相互作用は量子色力学（QCD）で記述され、その力を運ぶ粒子がグルーオンである。光子が電荷を持たないのと違い、グルーオンは自分も「色」の電荷を持ち、グルーオンどうしも結合する。そのためクォークとグルーオンは陽子の中に閉じ込められ、単独では取り出せない。

陽子の中を見る方法は、Rutherford が原子核を見つけたのと同じ、散乱である。1960年代の終わりに SLAC で、高いエネルギーの電子を陽子に当てると、大きな角度に散乱される電子が予想より多く、陽子の中に点状の粒子があるかのように振る舞った。Feynman はこの構成粒子を\textbf{パートン}（parton、part の粒子）と呼んだ。今では、パートンはクォーク、反クォーク、グルーオンのことである。

陽子の中には、$uud$ のほかに、グルーオンが量子的にクォークと反クォークの対（$g\to q\bar q$）に分かれてできる\textbf{海のクォーク}がある。だから陽子には $\bar u$、$\bar d$、$s$、$\bar s$、$c$、$b$ なども含まれる。ただし差し引きの数は決まっていて、$u$ の数から $\bar u$ の数を引くと2、$d-\bar d$ は1、$s-\bar s$ は0 になる（本ノートで使うパートン分布で、2.00、1.00、0.00）。この $uud$ の部分を\textbf{価クォーク}と呼ぶ。

\subsection{高速の陽子：パートン模型}\label{sec:partonmodel}

LHC では、陽子1個のエネルギーは 6.5~TeV で、Lorentz 因子は $\gamma=E/m_p\approx6900$ である。陽子の中の運動は時間の遅れで $\gamma$ 倍ゆっくりになり、陽子は進行方向に $1/\gamma$ に縮んだ薄い円盤になる。一方、パートンどうしが大きな運動量 $Q$ をやりとりする衝突は、時間 $\hbar/Q$（$Q=91$~GeV で $10^{-26}$~秒）で終わる。その短い間、パートンは互いの力を感じず、自由な粒子が陽子と同じ方向に並んで飛んでいるように見える。これが\textbf{パートン模型}である。

各パートンは陽子の運動量の割合 $x$（$0<x<1$）を持つ。\textbf{パートン分布関数}（PDF）$f_i(x)$ は、
\[
f_i(x)\,\dd x=\text{陽子の中で、種類 $i$ のパートンが運動量の割合 $x$ から $x+\dd x$ を持つ数の期待値}
\]
で定義される。全部のパートンの運動量を足すと陽子の運動量になるので、$\sum_i\int_0^1x f_i(x)\,\dd x=1$ である。本ノートで使う分布（NNPDF4.0、第~\ref{sec:setup} 節）では、$Q=M_Z$ でグルーオンが運動量の $47\%$、クォークと反クォークが $53\%$ を運ぶ。$x f(x)$ を図~\ref{fig:pdf} に示す。価クォーク（$u_v=u-\bar u$、$d_v=d-\bar d$）は $x\approx0.1$〜$0.3$ に山を持ち、海のクォークとグルーオンは小さい $x$ で急に大きくなる。

\begin{figure}[htbp]
\centering
\begin{tikzpicture}
\begin{semilogxaxis}[width=0.72\linewidth,height=6.2cm,xmin=1e-4,xmax=1,ymin=0,ymax=1.0,
  xlabel={$x$},ylabel={$x f(x)$},legend pos=north east,legend style={font=\small,draw=none,fill=none},
  tick label style={font=\small},label style={font=\small}]
\addplot[thick,blue] table[x=x,y=uv]{figs/pdf_mz.dat};\addlegendentry{$u_v=u-\bar u$}
\addplot[thick,red] table[x=x,y=dv]{figs/pdf_mz.dat};\addlegendentry{$d_v=d-\bar d$}
\addplot[thick,blue,dashed] table[x=x,y=ubar]{figs/pdf_mz.dat};\addlegendentry{$\bar u$}
\addplot[thick,red,dashed] table[x=x,y=dbar]{figs/pdf_mz.dat};\addlegendentry{$\bar d$}
\addplot[thick,green!50!black,dotted] table[x=x,y=s]{figs/pdf_mz.dat};\addlegendentry{$s$}
\addplot[thick,black] table[x=x,y=g10]{figs/pdf_mz.dat};\addlegendentry{$g/10$}
\end{semilogxaxis}
\end{tikzpicture}
\caption{本ノートで使うパートン分布（NNPDF4.0 NLO、$Q=M_Z$）の $x f(x)$。グルーオンは $\frac1{10}$ にして描いた。海のクォークとグルーオンは小さい $x$ で枠の外まで大きくなる（$x\bar u$ は $x<3\times10^{-3}$ で1を超え、$x=10^{-3}$ で $x\bar u=1.50$、$xg=29$）。$Z$ を $y=0$ で作るパートンは $x\approx0.007$ を持つ（第~\ref{sec:dylo} 節）。}
\label{fig:pdf}
\end{figure}

パートン分布は、陽子の中の強く結合した運動（摂動で計算できない部分）を表すので、計算では求めず、多くの実験のデータに合わせて決める（NNPDF などのグループが公開している）。同じ分布が、どの過程の計算にも使える。これが、パートン模型を予言の道具にしている点である。

二つの陽子（運動量 $P_1$、$P_2$、重心系のエネルギーの2乗 $S=(P_1+P_2)^2$、LHC で $\sqrt S=13$~TeV）から、パートン $i$（運動量 $x_1P_1$）と $j$（$x_2P_2$）が衝突するとする。陽子の質量を無視すると、パートンどうしの重心系のエネルギーの2乗は
\[
\hat s=(x_1P_1+x_2P_2)^2=2x_1x_2\,P_1\cdot P_2=x_1x_2S
\]
である（帽子は「パートンの」の印）。陽子どうしの断面積は、パートンの組ごとの断面積 $\hat\sigma_{ij}(\hat s)$ に、その組が衝突に来る確率 $f_i(x_1)\dd x_1\,f_j(x_2)\dd x_2$ を掛けて足したもの
\begin{equation}
\sigma(pp\to X)=\sum_{ij}\int_0^1\dd x_1\int_0^1\dd x_2\,f_i(x_1)f_j(x_2)\,\hat\sigma_{ij}(x_1x_2S)
\label{eq:partonmodel}
\end{equation}
になる。これがパートン模型の基本の式で、本ノートの式~\eqref{eq:had} はこれに QCD の補正を入れたものである。

\subsection{Drell--Yan 過程の最低次}\label{sec:dylo}

本ノートで扱う \textbf{Drell--Yan 過程}は、一方の陽子のクォークと他方の陽子の反クォークが対消滅して仮想の光子か $Z$ ボソン（電荷が合えば $W$ ボソン）になり、それがレプトンの対（$e^+e^-$ や $e^+\nu$）に崩壊する過程である（図~\ref{fig:dy}(a)）。$e^+e^-\to\mu^+\mu^-$ を逆向きにしたものと思ってよい。終状態がレプトンだけなので、測定がきれいで、パートン分布の検証と LHC の基準の測定に使われる。

\begin{figure}[htbp]
\centering
\begin{tikzpicture}[scale=0.82,every node/.style={font=\small}]
\begin{scope}
\node[blob] (b1) at (0,1.1) {};
\node[blob] (b2) at (0,-1.1) {};
\draw[thick,double] (-1.6,1.1) -- (b1) node[pos=0,left] {$p$};
\draw[thick,double] (-1.6,-1.1) -- (b2) node[pos=0,left] {$p$};
\foreach \dy in {0.05,0.2,0.35} {\draw[thin] (b1) -- ++(1.9,0.55+\dy); \draw[thin] (b2) -- ++(1.9,-0.55-\dy);}
\coordinate (v) at (1.7,0);
\draw[fermion] (b1) -- (v) node[midway,above right=-1pt] {$q\ (x_1)$};
\draw[fermion] (v) -- (b2) node[midway,below right=-1pt] {$\bar q\ (x_2)$};
\node[vtx] at (v) {};
\coordinate (w) at (3.2,0);
\draw[photon] (v) -- (w) node[midway,above] {$\gamma^*/Z$};
\node[vtx] at (w) {};
\draw[fermion] (4.3,0.85) -- (w) node[pos=0,right] {$\ell^+$};
\draw[fermion] (w) -- (4.3,-0.85) node[right] {$\ell^-$};
\node at (1.8,-2.3) {(a) LO};
\end{scope}
\begin{scope}[xshift=6.6cm]
\coordinate (a) at (0,1); \coordinate (b) at (0,-1); \coordinate (u1) at (1.2,0.55); \coordinate (u2) at (1.2,-0.55); \coordinate (v) at (1.9,0);
\draw[fermion] (a) -- (u1) node[pos=0,left] {$q$};
\draw[fermion] (u1) -- (v); \draw[fermion] (v) -- (u2); \draw[fermion] (u2) -- (b) node[left] {$\bar q$};
\draw[gluon] (u1) -- (u2) node[midway,left] {$g$};
\draw[photon] (v) -- (3.1,0) node[right] {$V^*$};
\foreach \p in {u1,u2,v} \node[vtx] at (\p) {};
\node at (1.4,-2.3) {(b) 仮想の補正};
\end{scope}
\begin{scope}[xshift=11.3cm]
\coordinate (a) at (0,1); \coordinate (b) at (0,-1); \coordinate (u1) at (0.9,0.55); \coordinate (v) at (1.6,0);
\draw[fermion] (a) -- (u1) node[pos=0,left] {$q$};
\draw[fermion] (u1) -- (v); \draw[fermion] (v) -- (b) node[left] {$\bar q$};
\draw[gluon] (u1) -- (2.4,1.25) node[right] {$g$};
\draw[photon] (v) -- (2.8,0) node[right] {$V^*$};
\foreach \p in {u1,v} \node[vtx] at (\p) {};
\node at (1.4,-2.3) {(c) 実の放射};
\end{scope}
\begin{scope}[xshift=15.6cm]
\coordinate (a) at (0,1); \coordinate (b) at (0,-1); \coordinate (u1) at (1.0,-0.45); \coordinate (v) at (1.0,0.45);
\draw[fermion] (a) -- (v) node[pos=0,left] {$q$};
\draw[fermion] (v) -- (u1);
\draw[gluon] (b) -- (u1) node[pos=0,left] {$g$};
\draw[fermion] (u1) -- (2.2,-1.1) node[right] {$q$};
\draw[photon] (v) -- (2.3,0.85) node[right] {$V^*$};
\foreach \p in {u1,v} \node[vtx] at (\p) {};
\node at (1.1,-2.3) {(d) $qg$ のチャンネル};
\end{scope}
\end{tikzpicture}
\caption{(a) Drell--Yan 過程の最低次（LO）。灰色の楕円は陽子で、残りの線は衝突に加わらなかったパートン。(b)〜(d) は NLO の QCD の補正の例。$V^*$ は仮想の $\gamma^*$、$Z$、$W$。(c) のグルーオンを放つ図と、反クォークから放つ図がある。(d) は一方の陽子のグルーオンがクォークと反クォークに分かれ、反クォークが他方のクォークと消滅する。}
\label{fig:dy}
\end{figure}

最低次（\textbf{LO}、leading order）では、レプトンの対の不変質量 $Q$ は $\sqrt{\hat s}$ に等しく、レプトンの対の運動量は $x_1P_1+x_2P_2$ である。レプトンの対の\textbf{ラピディティ} $y=\frac12\ln\frac{E+p_z}{E-p_z}$（ビームの方向へのブーストで定数だけずれる、速度に似た量）は $y=\frac12\ln(x_1/x_2)$ になる。だから $Q$ と $y$ を測ると
\[
x_1=\frac Q{\sqrt S}e^{y},\qquad x_2=\frac Q{\sqrt S}e^{-y}
\]
が決まる。$Z$（$Q=M_Z$）を $y=0$ で作るパートンは $x_1=x_2=0.0070$、$y=2.5$ では $x_1=0.085$、$x_2=0.00058$ で、ラピディティの分布はパートン分布の $x$ の依存をそのまま映す。$Z$ の質量の付近（66〜116~GeV）の LO の断面積の内訳は、$u\bar u$ が $39\%$、$d\bar d$ が $37\%$、$s\bar s$ が $15\%$、$c\bar c$ が $6\%$、$b\bar b$ が $3\%$ である。反クォークはすべて海から来るので、海のクォークがなければ Drell--Yan はほとんど起こらない。$W^+$ は主に $u\bar d$ から、$W^-$ は $d\bar u$ から作られる。陽子の中の価クォークは $u$ が $d$ の2倍あるので $W^+$ のほうが多く、比は $\sigma(W^+)/\sigma(W^-)=1.34$ になる（本ノートの NLO の値。海のクォークの寄与があるので2にはならない）。

パートンどうしの断面積 $\hat\sigma$ は、学部の電磁気の延長で、$q\bar q\to\gamma^*\to e^+e^-$ なら $\hat\sigma=\frac{4\pi\alpha^2Q_q^2}{3N_c\hat s}$（$e^+e^-\to\mu^+\mu^-$ の $\frac{4\pi\alpha^2}{3s}$ に、クォークの電荷 $Q_q$ の2乗と、色の平均 $1/N_c=\frac13$ を掛けたもの）である。$Z$ を交換する部分は、$Z$ の質量で共鳴（Breit--Wigner の形）を作る（式~\eqref{eq:born}）。断面積の単位は pb（ピコバーン、$1~\mathrm{pb}=10^{-36}~\mathrm{cm}^2$）で、$Z\to e^+e^-$ の約 2~nb $=2000$~pb は、LHC の Run 2 のデータ量（積分ルミノシティ $140~\mathrm{fb}^{-1}$）で $2.8\times10^8$ 事象に当たる。

\subsection{強い相互作用の補正：なぜ NLO が要るか}\label{sec:whynlo}

式~\eqref{eq:partonmodel} の $\hat\sigma$ は、強い相互作用の結合定数 $\alpha_s$（電磁気の $\alpha=\frac1{137}$ にあたる）の冪級数で計算する。$\alpha_s$ を含まない項が LO、$\alpha_s$ の1次の補正を加えたものが \textbf{NLO}（next-to-leading order）である。$\alpha_s$ は量子補正のためにエネルギーとともに小さくなる（\textbf{漸近的自由}）。本ノートの値で $\alpha_s=0.297$（2~GeV）、$0.118$（$M_Z$）、$0.087$（1~TeV）で、近似的に
\[
\alpha_s(\mu)\approx\frac{\alpha_s(M_Z)}{1+b_0\,\alpha_s(M_Z)\ln(\mu^2/M_Z^2)},\qquad b_0=\frac{33-2n_f}{12\pi}
\]
と振る舞う（$n_f=5$ はクォークの種類の数）。$\alpha_s(M_Z)/\pi\approx0.04$ は小さいが、Drell--Yan の NLO の補正は断面積を約 $14\%$ 上げる（$K$ 因子 $\approx1.14$）。LHC の測定の精度（数 $\%$）で比べるには NLO 以上が要る。

NLO の補正には三種類ある（図~\ref{fig:dy}(b)〜(d)）。(b) \textbf{仮想の補正}：クォークと反クォークがグルーオンを交換してから消滅する（ループの図）。終状態は LO と同じである。(c) \textbf{実の放射}：消滅の前後にグルーオンが1個放たれる。(d) 一方の陽子の\textbf{グルーオンが始状態}になる $qg\to V^*q$。電磁気で、加速された電荷が光子を放つ制動放射（bremsstrahlung）と同じで、(c) は QCD の制動放射である。

\subsection{発散と、それがパートン分布に移るわけ}\label{sec:div}

(b) と (c) は、そのままでは無限大になる。理由は電磁気の制動放射と同じである。電荷が光子を放つ確率は、光子のエネルギー $\omega$ について $\dd\omega/\omega$ に比例し、$\omega\to0$ で積分が対数的に発散する（\textbf{軟らかい発散}）。さらに、質量0のクォークがグルーオンを放つときは、放つ角度 $\theta$ についても $\dd\theta^2/\theta^2$ の形になり、グルーオンがクォークと平行（$\theta\to0$）なところで発散する（\textbf{共線の発散}）。式~\eqref{eq:real} の $1/t$、$1/u$（$t\propto\theta^2$）がこれである。

軟らかい発散は、実の放射と仮想の補正の和で消える。エネルギー0のグルーオンを放った事象は、放たなかった事象と区別できないので、両方を足したものだけが測れる量であり、その和は有限になる（Bloch--Nordsieck の定理。崩壊の NLO のノートでも同じことを確かめた）。

始状態の共線の発散は消えない。しかし物理の意味は明らかである。陽子の中で運動量の割合 $x$ のクォークが、ほぼ平行なグルーオンを放って割合 $zx$ になってから衝突したとする。これは、初めから割合 $zx$ のクォークが衝突したのと区別できない。つまりこの発散は、衝突そのものではなく\textbf{陽子の中の構造}の一部で、パートン分布がすでに表しているものを、もう一度数えている。そこで、平行な放射の部分を $\hat\sigma$ から差し引き、パートン分布に入れる。これを\textbf{因子化}（factorization）と呼ぶ。分け目の運動量（おおまかには、放射の横向きの運動量が $\mu_F$ より小さいものはパートン分布へ、大きいものは $\hat\sigma$ へ。実際の分け方は次の $\MSbar$ 方式で決める）を\textbf{因子化のスケール} $\mu_F$ と呼ぶ。

計算の途中で無限大を扱うため、時空の次元を $d=4-2\epsilon$ に変える\textbf{次元正則化}を使う。たとえば $\int_0^1\dd t\,t^{-1-\epsilon}=-1/\epsilon$（$\epsilon<0$ で収束する積分を解析接続したもの）のように、対数の発散が $1/\epsilon$ の極になる。物理の量では、実と仮想の和と、パートン分布への移し替えで極がすべて消え、最後に $\epsilon\to0$ とする（第~\ref{sec:coef} 節、検算の手順 (1)(d)）。

移し替えるときに、極と一緒に有限の部分をどれだけ移すかには任意性がある。標準の約束が\textbf{$\MSbar$（エムエスバー）方式}で、極（と決まった定数 $\ln4\pi-\gamma_E$）だけを移す。データから決めたパートン分布は $\MSbar$ 方式で定義されているので、$\hat\sigma$ の側も $\MSbar$ 方式で計算しなければならない。方式がずれると、有限部分がずれた断面積になる。

\subsection{二つのスケールとパートン分布の進化}\label{sec:scales}

パートン分布は $\mu_F$ の関数 $f_i(x,\mu_F)$ になる。$\mu_F$ を上げると、より細かく陽子を見ることになり、クォークがグルーオンを放っていた様子や、グルーオンがクォークの対に分かれていた様子が見えてくる。その変化は \textbf{DGLAP の方程式}
\[
\frac{\partial f_q(x,\mu_F)}{\partial\ln\mu_F^2}=\frac{\alpha_s}{2\pi}\int_x^1\frac{\dd z}z\Bigl[P_{qq}(z)\,f_q\Bigl(\frac xz,\mu_F\Bigr)+P_{qg}(z)\,f_g\Bigl(\frac xz,\mu_F\Bigr)\Bigr]
\]
で決まる。\textbf{分岐関数} $P_{qq}(z)=C_F\frac{1+z^2}{1-z}$（$z<1$）は、クォークがグルーオンを放って運動量の割合 $z$ を残す確率の密度、$P_{qg}(z)=T_R[z^2+(1-z)^2]$ はグルーオンがクォークの対に分かれ、クォークが割合 $z$ を持つ確率の密度である（$C_F=\frac43$、$T_R=\frac12$ は色の数え上げの因子）。本ノートの分布では、$x=0.3$ の $xu$ は $Q=10$、100、1000~GeV で 0.453、0.339、0.274 と減り、$x=0.001$ の $xg$ は 16.2、29.7、39.0 と増える。大きい $x$ のパートンが放射で小さい $x$ に移るためである。実は、第~\ref{sec:div} 節の共線の極の係数は、ちょうどこの分岐関数になる（式~\eqref{eq:dqq} の $L_FP_{qq}$）。

もう一つのスケールは、$\alpha_s(\mu_R)$ を測る\textbf{くりこみのスケール} $\mu_R$ である。本当の断面積は $\mu_F$ にも $\mu_R$ にもよらないが、摂動の級数を途中で切ると、その次の次数の大きさだけ依存が残る。そこで $\mu_F$、$\mu_R$ を典型的な値（Drell--Yan では $Q$ や $M_Z$）の $\frac12$ 倍から2倍に動かした変化を、計算していない高次の項の大きさの目安にする。本ノートでは、$Z$ の断面積のこの変化が LO の $\pm12\%$ から NLO の $^{+4}_{-7}\%$ に縮む（表~\ref{tab:dyk}）。

\subsection{係数関数とプラスの分布}\label{sec:coefintro}

NLO の $\hat\sigma$ から発散を移し替えた残りを、LO の断面積 $\sigma_0$ で割ったものが\textbf{係数関数} $\Delta_{ij}(z)$ である（式~\eqref{eq:had}）。$z=Q^2/\hat s$ は、パートンの衝突のエネルギーのうちレプトンの対に行った割合で、$1-z$ は放たれたグルーオン（またはクォーク）が持ち去った割合を表す。LO ではすべてがレプトンの対に行くので $\Delta=\delta(1-z)$（デルタ関数）である。

実の放射は $z<1$ の滑らかな関数だが、$z\to1$（放射が軟らかい）で $1/(1-z)$ のように大きくなり、仮想の補正はちょうど $z=1$ にだけある（$\delta(1-z)$ に比例し、係数が無限大）。二つを合わせた有限のものが\textbf{プラスの分布}
\[
\int_0^1\dd z\Bigl[\frac1{1-z}\Bigr]_+h(z)=\int_0^1\dd z\,\frac{h(z)-h(1)}{1-z}
\]
で、なめらかな関数 $h$ との積分の形で定義される（$h(1)$ を引くことで $z\to1$ の発散が打ち消される）。式~\eqref{eq:dqq} の $[\ln(1-z)/(1-z)]_+$ も同じ形である。陽子の断面積は、係数関数とパートン分布を掛けて $x_1$、$x_2$（または $z$）で積分した式~\eqref{eq:lumi} で求まる。NLO で LO から増えた割合を \textbf{$K$ 因子}と呼ぶ。$Z$ の質量の付近では、$q\bar q$ のチャンネルが $+29\%$、$qg$ のチャンネルが $-15\%$ で、$K=1.14$ になる（表~\ref{tab:dyk}）。$qg$ が負になるのは、移し替えた共線の部分を引いた残りだからで、確率として負になったわけではない。

\subsection{ヒッグスとグルーオン：重いトップの有効理論}\label{sec:heftintro}

LHC でヒッグスボソンが作られる主な過程は、二つのグルーオンの融合 $gg\to H$ である。ヒッグスは粒子の質量に比例して結合するので、質量0のグルーオンには直接結合しない。しかし、一番重いトップクォークのループを通して結合できる（図~\ref{fig:ggh}(a)）。この過程は最低次で $\alpha_s^2$ に比例する。トップの質量 $m_t=172.5$~GeV はヒッグスの質量の半分 $M_H/2=62.6$~GeV よりずっと大きいので、ループは点のように見え、その効果は次の\textbf{有効ラグランジアン}にまとめられる（図~\ref{fig:ggh}(b)）。
\[
\mathcal L_{\mathrm{eff}}=-\frac14c_{Hgg}\frac Hv\,G^a_{\mu\nu}G^{a\mu\nu},\qquad c_{Hgg}=-\frac{\alpha_s}{3\pi}\Bigl(1+\frac{11}4\frac{\alpha_s}\pi+\cdots\Bigr)
\]
（$v=246$~GeV はヒッグスの真空期待値、$G^a_{\mu\nu}$ はグルーオンの場の強さ。括弧の中の $\frac{11}4\frac{\alpha_s}\pi$ は、トップのループ自身の QCD の補正で、\textbf{ウィルソン係数}の補正と呼ぶ）。これが\textbf{重いトップの有効理論}で、ヒッグスのノート~\cite{higgsnote} の $H\to gg$ と同じものである。トップの質量を正しく入れた LO の断面積は、この極限より $6.6\%$ 大きい（本ノートの計算で $|A_{1/2}(\tau_t)/A_{1/2}(0)|^2=1.066$、$\tau_t=M_H^2/4m_t^2$）。

\begin{figure}[htbp]
\centering
\begin{tikzpicture}[scale=0.85,every node/.style={font=\small}]
\begin{scope}
\coordinate (a) at (0,1); \coordinate (b) at (0,-1); \coordinate (t1) at (1.3,0.7); \coordinate (t2) at (1.3,-0.7); \coordinate (t3) at (2.4,0);
\draw[gluon] (a) -- (t1) node[pos=0,left] {$g$};
\draw[gluon] (b) -- (t2) node[pos=0,left] {$g$};
\draw[fermion] (t1) -- (t3); \draw[fermion] (t3) -- (t2); \draw[fermion] (t2) -- (t1) node[midway,left] {$t$};
\draw[thick,dashed] (t3) -- (3.6,0) node[right] {$H$};
\foreach \p in {t1,t2,t3} \node[vtx] at (\p) {};
\node at (1.6,-1.7) {(a) トップのループ};
\end{scope}
\begin{scope}[xshift=6cm]
\coordinate (a) at (0,1); \coordinate (b) at (0,-1); \coordinate (v) at (1.4,0);
\draw[gluon] (a) -- (v) node[pos=0,left] {$g$};
\draw[gluon] (b) -- (v) node[pos=0,left] {$g$};
\draw[thick,dashed] (v) -- (2.6,0) node[right] {$H$};
\node[rectangle,fill=black,inner sep=2.5pt] at (v) {};
\node at (1.2,-1.7) {(b) 有効な頂点（$m_t\to\infty$）};
\end{scope}
\begin{scope}[xshift=11cm]
\coordinate (a) at (0,1); \coordinate (b) at (0,-1); \coordinate (v) at (1.4,0); \coordinate (u) at (0.7,0.5);
\draw[gluon] (a) -- (u) node[pos=0,left] {$g$};
\draw[gluon] (u) -- (v);
\draw[gluon] (u) -- (1.6,1.3) node[right] {$g$};
\draw[gluon] (b) -- (v) node[pos=0,left] {$g$};
\draw[thick,dashed] (v) -- (2.6,0) node[right] {$H$};
\node[rectangle,fill=black,inner sep=2.5pt] at (v) {};
\node[vtx] at (u) {};
\node at (1.2,-1.7) {(c) 実の放射の例};
\end{scope}
\end{tikzpicture}
\caption{(a) $gg\to H$ の最低次（トップのループ）。(b) 重いトップの極限で点になった頂点（黒い四角）。(c) NLO の実の放射の一つ。有効理論では、ループがすでに一つ縮んでいるので、NLO の補正は Drell--Yan と同じ1ループと実の放射の計算になる。}
\label{fig:ggh}
\end{figure}

NLO の補正の考え方は Drell--Yan と同じで、仮想の補正、グルーオンを1個放つ実の放射（$gg\to Hg$）、新しいチャンネル（$qg\to Hq$、$q\bar q\to Hg$）があり、始状態の共線の発散をパートン分布に移す。違いは二つある。一つは、始状態がグルーオンなので、グルーオンが放射する強さを表す色の因子が $C_A=3$ になり（クォークは $C_F=\frac43$）、補正が大きいことである。本ノートの値で、$\mu=M_H$ の NLO は LO の2.29倍になる（$K=2.29$）。もう一つは、LO が $\alpha_s^2$ に比例するので、スケールの依存が LO から大きい（$\mu_R=\mu_F$ を2倍、半分にして $-14\%$、$+17\%$）ことである。

\subsection{\MGaMC{} がすること}\label{sec:mg5intro}

本ノートの検算の相手 \MGaMC{} は、過程を指定すると Feynman 図を自動で作り、振幅を数値で計算し、位相空間の積分をモンテカルロ法（乱数で点を選んで平均する）で行うプログラムである。そのため結果は統計の誤差を持つ。NLO（\texttt{[QCD]} と指定する）では、実の放射の発散を、同じ発散を持つ簡単な関数を引いて数値で積分できる形にし（FKS の引き算~\cite{fks}）、引いた分を解析的に足し戻す。本ノートの係数関数は解析的に求めた独立の結果なので、両者が一致すれば、\MGaMC{} の NLO の仕組み（引き算、仮想の補正、パートン分布との畳み込み、スケールを変えた重みの付け直し）がこの過程で正しく働いていることになる。手順 (3) の MC@NLO は、NLO の計算に、パートンが次々に放射を繰り返す様子（パートンシャワー）を数え過ぎないようにつなぐ方法で、実際の事象の模擬に使う。

\begin{table}[htbp]
\centering\small
\caption{本ノートの用語と記号。}
\label{tab:glossary}
\begin{tabularx}{\linewidth}{>{\raggedright\arraybackslash}p{3.6cm}X}
\toprule
\textbf{用語、記号} & \textbf{意味} \\
\midrule
パートン & 陽子の構成粒子（クォーク、反クォーク、グルーオン） \\
パートン分布（PDF）$f_i(x,\mu_F)$ & パートン $i$ が運動量の割合 $x$ を持つ数の密度（第~\ref{sec:partonmodel} 節） \\
$S$、$\hat s=x_1x_2S$ & 陽子どうし、パートンどうしの重心系のエネルギーの2乗 \\
$Q$、$\tau=Q^2/S$、$z=Q^2/\hat s$ & レプトンの対の不変質量、その2乗の $S$、$\hat s$ に対する割合 \\
ラピディティ $y$ & $\frac12\ln\frac{E+p_z}{E-p_z}$。LO で $\frac12\ln(x_1/x_2)$ \\
LO、NLO & 摂動の最低次、$\alpha_s$ の1次まで \\
$\alpha_s(\mu_R)$、$a_s=\alpha_s/4\pi$ & 強い相互作用の結合定数。$\mu_R$ はそれを測るスケール \\
$\mu_F$ & 因子化のスケール（パートン分布と $\hat\sigma$ の分け目） \\
仮想の補正、実の放射 & ループの図による補正、グルーオンを1個放つ過程 \\
軟らかい発散、共線の発散 & 放射のエネルギー $\to0$、放射の角度 $\to0$ で起こる無限大 \\
次元正則化、$\epsilon$ & 時空の次元を $4-2\epsilon$ にして発散を $1/\epsilon$ の極にする方法 \\
因子化、$\MSbar$ 方式 & 共線の発散をパートン分布に移すこと、その標準の約束 \\
DGLAP の方程式、分岐関数 $P_{ij}$ & パートン分布の $\mu_F$ による変化、その変化の確率の密度 \\
係数関数 $\Delta_{ij}(z)$ & 発散を移したあとのパートンの断面積を $\sigma_0$ で割ったもの \\
プラスの分布 $[\,\cdot\,]_+$ & $z\to1$ の軟らかい発散を仮想の補正と打ち消した分布 \\
$K$ 因子 & NLO の断面積と LO の断面積の比 \\
$C_F=\frac43$、$T_R=\frac12$、$N_c=3$ & QCD の色の数え上げの因子 \\
$V^*$、$\sigma_0$ & 仮想の $\gamma^*$、$Z$、$W$、LO のパートンの断面積 \\
pb、fb & 断面積の単位（$10^{-36}$、$10^{-39}~\mathrm{cm}^2$） \\
固定次数、MC@NLO & シャワーなしの NLO の計算、シャワーとつないだ NLO \\
パートンシャワー & パートンが次々に放射する様子を確率的に作るプログラム（Pythia 8 など。第~\ref{sec:step3} 節） \\
Sudakov の因子 & ある $p_T$ より上では放射しない確率 \\
プル & 2つの値の差を誤差で割ったもの（統計の一致の目安） \\
重いトップの有効理論 & トップのループを点にした $ggH$ の結合（第~\ref{sec:heftintro} 節） \\
ウィルソン係数 & 有効な結合の係数 $c_{Hgg}$ とその $\alpha_s$ の補正 $\frac{11}4\frac{\alpha_s}\pi$ \\
紫外の対抗項 & ループの紫外の発散を打ち消す頂点（結合定数の繰り込みから来る） \\
R2 の頂点 & ループを4次元の運動量で組み立てるときに落ちる有限な項を補う頂点（第~\ref{sec:step2} 節 (i)） \\
UFO & \MGaMC{} が読む模型のファイル（粒子、頂点、NLO では紫外の対抗項と R2 の頂点） \\
$C_A=3$ & グルーオンが放射する強さの色の因子 \\
\bottomrule
\end{tabularx}
\end{table}

\section{設定}\label{sec:setup}

\paragraph{模型と過程。}クォークは質量0の5フレーバー（$u$、$d$、$s$、$c$、$b$。$b$ は陽子の中にある）、CKM 行列は単位行列とする。色の因子は $C_F=\frac43$、$T_R=\frac12$、$N_c=3$ である。$d$ 次元の計算では、ベクトルの流れ $\gamma^\mu$ と軸性の流れ $\gamma^\mu\gamma_5$（$\gamma_5$ は反可換）を手で書き、4次元の振幅の2乗を NLO ノートの FeynArts の模型（\texttt{SM\_CKM\_FA}、本ノートの模型と同じもの）と比べる。\MGaMC{} 3.8.1 では、同梱の \texttt{loop\_sm} を $b$ の質量0の制限（\texttt{loop\_sm-no\_b\_mass}）で使う。これがツールとして確かめる対象である。

\paragraph{入力。}電弱のパラメータは \texttt{loop\_sm} の param\_card の値とし、両方の計算で同じものを使う：$\alpha^{-1}=132.507$、$G_F=1.16639\times10^{-5}$~GeV$^{-2}$、$M_Z=91.188$~GeV、そこから $M_W=80.419002$~GeV、$s_W^2=1-M_W^2/M_Z^2=0.222246$。$V$ の伝播関数は固定の幅の Breit--Wigner で、$\Gamma_Z=2.441404$~GeV、$\Gamma_W=2.0476$~GeV である。これらは木の関係で決まる値で、SM ノートの $\MSbar$ のパラメータではない。本ノートの比べる量は、両方の計算が同じ入力から求めた断面積である。パートン分布は NNPDF4.0 NLO~\cite{nnpdf40}（LHAPDF~6.5.5 の \texttt{NNPDF40\_nlo\_as\_01180}、番号 331700）で、$\alpha_s(\mu_R)$ もこの組のもの（2ループの走り、$\alpha_s(M_Z)=0.118$、$n_f\le5$）を使う。$\sqrt S=13$~TeV、$\mu_R$ と $\mu_F$ は固定の値 $\mu_0$ とその $\frac12$、2倍である。

\paragraph{断面積の定義。}レプトンの対の不変質量を $Q$、$\tau=Q^2/S$、$z=Q^2/\hat s$ とする。レプトンの角度は積分するので、レプトンのテンソルは $(q^\mu q^\nu-Q^2g^{\mu\nu})$ に $Q^2$ の関数を掛けたものになり、流れが保存するので $-g_{\mu\nu}H^{\mu\nu}$ だけが効く（$H^{\mu\nu}$ はクォークの側の振幅の2乗を、$V^*$ の偏極の添字をつけたまま書いたもの。レプトンの飛ぶ向きを積分した断面積は、そのうち $V^*$ の偏極を平均した一つの組み合わせだけで決まる）。
\begin{equation}
\begin{gathered}
\frac{\dd\sigma}{\dd Q^2}=\sum_{ij}\int\dd x_1\dd x_2\,f_i(x_1,\mu_F)f_j(x_2,\mu_F)\,\frac{\sigma_0(Q^2)}{\hat s}\,\Delta_{ij}\Bigl(z,\frac{Q^2}{\mu_F^2}\Bigr),\\
\Delta_{ij}=\delta_{ij}^{(0)}\delta(1-z)+a_s\Delta^{(1)}_{ij}+\cdots,\qquad a_s=\frac{\alpha_s(\mu_R)}{4\pi}.
\end{gathered}
\label{eq:had}
\end{equation}
$\sigma_0(Q^2)$ は $q\bar q'\to V^*\to$ レプトンの Born の断面積の $\hat s=Q^2$ での値で、$qg$ のチャンネルはそのクォークのフレーバーの $\sigma_0$ を掛ける。

\section{結果}\label{sec:results}

この節は、第~\ref{sec:primer} 節の考え方を式にしたものである。第~\ref{sec:born} 節は LO の $\sigma_0$ と、パートン模型の式~\eqref{eq:partonmodel} を $z$ の積分に書き直した形（第~\ref{sec:dylo} 節）、第~\ref{sec:coef} 節は NLO の三種類の補正（図~\ref{fig:dy}(b)〜(d)）と、発散をパートン分布に移したあとの係数関数（第~\ref{sec:div}〜\ref{sec:coefintro} 節）である。$t$、$u$ は放射の角度を、$v$ はその角度を $0$ から $1$ の変数にしたものを表す。

\subsection{Born の断面積とハドロンの式}\label{sec:born}

$q\bar q\to\gamma^*/Z\to e^+e^-$ の Born は、クォークとレプトンのヘリシティ $i,j\in\{L,R\}$ について
\begin{equation}
\sigma_0(Q^2)=\frac1{N_c}\frac{Q^2}{48\pi}\sum_{i,j}|A_{ij}|^2,\qquad
A_{ij}=\frac{e^2Q_qQ_e}{Q^2}+\frac{e^2}{s_W^2c_W^2}\frac{g_i^qg_j^e}{Q^2-M_Z^2+iM_Z\Gamma_Z},\qquad g_L=T_3-Qs_W^2,\ g_R=-Qs_W^2
\label{eq:born}
\end{equation}
で、$u\bar d\to W^+\to e^+\nu$ は $A_{LL}=\frac{g^2}2(Q^2-M_W^2+iM_W\Gamma_W)^{-1}$（$g=e/s_W$）だけである。$\hat s=x_1x_2S$ から $(x_1,z)$ に変数を替えると、式~\eqref{eq:had} は
\begin{equation}
\frac{\dd\sigma}{\dd Q^2}=\frac{\sigma_0(Q^2)}{Q^2}\sum_{ij}\int_\tau^1\dd z\,\Delta_{ij}(z)\,G_{ij}\Bigl(\frac\tau z\Bigr),\qquad
G_{ij}(y)=\int_{\ln y}^0\dd t\,x f_i(e^t)\,x f_j(ye^{-t})
\label{eq:lumi}
\end{equation}
になる（$G=yL$、$L$ はパートンの光度）。プラスの分布は $\int\dd z\,[g]_+G(\tau/z)=\int_\tau^1\dd z\,g(z)[G(\tau/z)-G(\tau)]-G(\tau)\int_0^\tau g$ で扱う。

\subsection{NLO の係数関数}\label{sec:coef}

\paragraph{実の放射。}$q(p_1)\bar q(p_2)\to V^*(q)g(k)$ と $q(p_1)g(p_2)\to V^*(q)q(k)$ で、$t=(p_1-k)^2=-\hat s(1-z)v$、$u=(p_2-k)^2=-\hat s(1-z)(1-v)$ とする。$-g_{\mu\nu}H^{\mu\nu}$ の Born に対する比に、色と平均の因子（$q\bar q$ は $C_F$、$qg$ は $T_R/(1-\epsilon)$：グルーオンの偏極は $d-2$ 個）を掛けたものは
\begin{equation}
r_{q\bar q}=2C_F\Bigl[(1-\epsilon)\Bigl(\frac ut+\frac tu\Bigr)+\frac{2Q^2\hat s}{ut}-2\epsilon\Bigr],\qquad
r_{qg}=-\frac{2T_R}{1-\epsilon}\Bigl[(1-\epsilon)\Bigl(\frac u{\hat s}+\frac{\hat s}u\Bigr)+\frac{2Q^2t}{u\hat s}-2\epsilon\Bigr]
\label{eq:real}
\end{equation}
である（ベクトルと軸性の流れで同じ）。裸の結合 $g_0^2\mu_0^{2\epsilon}=4\pi\alpha_s\mu_R^{2\epsilon}(e^{\gamma_E}/4\pi)^\epsilon$ と2体の位相空間 $\int\dd\Phi_2=\frac1{8\pi}(\frac{4\pi}{\hat s})^\epsilon\frac{(1-z)^{1-2\epsilon}}{\Gamma(1-\epsilon)}\int_0^1\dd v\,[v(1-v)]^{-\epsilon}$ から
\begin{equation}
\frac{\Delta_R}{a_s}=\frac{e^{\epsilon\gamma_E}}{\Gamma(1-\epsilon)}\Bigl(\frac{\mu_R^2}{Q^2}\Bigr)^\epsilon z^\epsilon(1-z)^{1-2\epsilon}\int_0^1\dd v\,[v(1-v)]^{-\epsilon}\,r(z,v)
\label{eq:dr}
\end{equation}
で、$v$ の積分はベータ関数になる。$q\bar q$ では $(1-z)^{-1-2\epsilon}=-\frac1{2\epsilon}\delta(1-z)+[\frac1{1-z}]_++\cdots$ と展開する。

\paragraph{仮想の補正。}Feynman ゲージの頂点の補正を TID で $B_0(Q^2)$ と $C_0$ に直し、IBP の関係 $C_0=-\frac{2(d-3)}{(d-4)Q^2}B_0$ を使うと、時間的な $Q^2$ の形状因子は
\begin{equation}
F^{(1)}=-a_sC_F\Bigl(\frac{\mu_R^2}{-Q^2-i0}\Bigr)^\epsilon\frac{e^{\epsilon\gamma_E}\Gamma(1+\epsilon)\Gamma(1-\epsilon)^2}{\Gamma(1-2\epsilon)}\,\frac{2-\epsilon+2\epsilon^2}{\epsilon^2(1-2\epsilon)}
=-a_sC_F\Bigl(\frac{\mu_R^2}{-Q^2-i0}\Bigr)^\epsilon\Bigl[\frac2{\epsilon^2}+\frac3\epsilon+8-\zeta_2+O(\epsilon)\Bigr]
\label{eq:ff}
\end{equation}
で、$\Delta_V=2\re F^{(1)}\delta(1-z)$ である。質量0のクォークの自己エネルギーはスケールを持たないので0で、場の定数は効かない。$(-1-i0)^{-\epsilon}=e^{i\pi\epsilon}$ の $\cos\pi\epsilon$ が $-\frac{\pi^2}2\times$ 2重極、つまり $\delta(1-z)$ の係数に $+12\zeta_2C_F$ を足す。

\paragraph{共線の対抗項と結果。}$\MSbar$ のパートン分布に移す対抗項は $\frac{\alpha_s}{2\pi}\frac1\epsilon(\frac{\mu_R^2}{\mu_F^2})^\epsilon$ に $2P_{qq}$（$q\bar q$、2本の脚）、$P_{qg}$（$qg$）を掛けたもので、$P_{qq}=C_F[\frac{1+z^2}{1-z}]_+$、$P_{qg}=T_R[z^2+(1-z)^2]$ である。$\Delta_V+\Delta_R$ に足すと極が消え、$\ln\mu_R$ が落ちて（$L_F=\ln(\mu_F^2/Q^2)$）
\begin{align}
\Delta^{(1)}_{q\bar q}&=C_F\Bigl\{16\Bigl[\frac{\ln(1-z)}{1-z}\Bigr]_+-8(1+z)\ln(1-z)-4\frac{1+z^2}{1-z}\ln z+(8\zeta_2-16)\delta(1-z)\Bigr\}-4L_F\,P_{qq}(z),\label{eq:dqq}\\
\Delta^{(1)}_{qg}&=T_R\Bigl\{2[z^2+(1-z)^2]\ln\frac{(1-z)^2}z+1+6z-7z^2\Bigr\}-2L_F\,P_{qg}(z)\label{eq:dqg}
\end{align}
になる。$\delta(1-z)$ の係数 $8\zeta_2-16$ は、実の放射の $-6\zeta_2$ と仮想の補正の $-16+14\zeta_2$ の和である。仮想の $14\zeta_2$ のうち $12\zeta_2$ は時間的な $Q^2$ への解析接続（$\cos\pi\epsilon$）から来る。$\mu_F$ の依存 $-4L_FP_{qq}$、$-2L_FP_{qg}$ は、LO の $f(\mu_F)$ の DGLAP の進化を打ち消す。

\subsection{$gg\to H$ の NLO の係数関数（重いトップの有効理論）}\label{sec:gghcoef}

Drell--Yan と同じく、パートンの断面積を $\hat\sigma_{ij}=\sigma_0z\,\Delta_{ij}(z)$、$z=M_H^2/\hat s$ と書く（文献の係数関数は $z\Delta$ である）。LO は
\begin{equation}
\sigma_0=\frac{\alpha_s(\mu_R)^2}{576\pi v^2}\quad(c_{Hgg}=-\alpha_s/3\pi),\qquad
\sigma=\sigma_0\sum_{ij}\int_\tau^1\dd z\,\Delta_{ij}(z)\,G_{ij}\Bigl(\frac\tau z\Bigr),\quad\tau=\frac{M_H^2}S
\label{eq:ggh0}
\end{equation}
である。NLO では、$gg\to Hg$（$s$、$t$、$u$ チャンネルのグルーオンと $gggH$ の接触の頂点）、$qg\to Hq$、$q\bar q\to Hg$ の実の放射と、有効な頂点の1ループの補正（グルーオンとゴーストのループ）を $d$ 次元で求める。$c_{Hgg}\propto\alpha_s$ なので、紫外の対抗項は $\MSbar$ の $2\delta\alpha_s/\alpha_s=-2\beta_0a_s/\epsilon$（$\beta_0=11-\frac23n_f$、$\alpha_sG^2$ はこの次数で有限）である。共線の対抗項には $P_{gg}(z)=2C_A[\frac z{(1-z)_+}+\frac{1-z}z+z(1-z)]+\frac{\beta_0}2\delta(1-z)$ と $P_{gq}(z)=C_F\frac{1+(1-z)^2}z$ が入る。結果は（$L_F=\ln(\mu_F^2/M_H^2)$、$L_R=\ln(\mu_R^2/M_H^2)$、単位は $a_s$）
\begin{align}
\Delta^{(1)}_{gg}&=(24\zeta_2+2\beta_0L_R-2\beta_0L_F)\delta(1-z)+48\Bigl[\frac{\ln(1-z)}{1-z}\Bigr]_+-24L_F\Bigl[\frac1{1-z}\Bigr]_++24L_F\Bigl(2-\frac1z-z+z^2\Bigr)\notag\\
&\quad+\frac{22(1-z)^4-48(1-3z+3z^2-2z^3+z^4)\ln(1-z)+24(1-z+z^2)^2\ln z}{(z-1)z},\label{eq:dgg}\\
\Delta^{(1)}_{qg}&=\frac4{3z}\bigl[-3+6z-z^2+2(2-2z+z^2)(2\ln(1-z)-\ln z)\bigr]-\frac{8L_F}{3z}(2-2z+z^2),\qquad
\Delta^{(1)}_{q\bar q}=\frac{128(1-z)^3}{27z}\label{eq:dqgh}
\end{align}
で、ウィルソン係数の補正は $\Delta^{(1)}_{gg}$ の $\delta(1-z)$ に $4\times\frac{11}2=22$ を足す。$z\Delta^{(1)}$ は Dawson、Djouadi--Spira--Zerwas、Harlander--Kilgore の形~\cite{dawson91,dsz91,hk02}（参照）と記号として一致する（検算の手順 (2)）。$\delta(1-z)$ の $24\zeta_2=4\pi^2$ は、Drell--Yan の $8\zeta_2C_F$ と同じく時間的な形状因子の $\pi^2$ から来て、$C_F$ が $C_A$ に替わって大きくなる。

\section{検算の計画}\label{sec:plan}

各量を、どの手段で計算して突き合わせるかを表~\ref{tab:plan} に示す。手計算を基準とし、ツールの出力がそれと一致するかを確かめる。ツールの版と配置は計算環境の記録~\cite{env} にある。文献の値は参照として並べ、比べる量は本ノートの式と入力で求めたものにする。

\begin{table}[htbp]
\centering\small
\caption{各量と、それを計算して突き合わせる手段。}
\label{tab:plan}
\begin{tabularx}{\linewidth}{>{\raggedright\arraybackslash}p{3.4cm}>{\raggedright\arraybackslash}p{1.6cm}>{\raggedright\arraybackslash}X>{\raggedright\arraybackslash}p{3.6cm}}
\toprule
\textbf{量} & \textbf{手計算} & \textbf{FeynCalc + FeynArts、自分の積分} & \textbf{\MGaMC{}} \\
\midrule
Drell--Yan の係数関数 \eqref{eq:dqq}、\eqref{eq:dqg} & 第~\ref{sec:coef} 節 & $d$ 次元の実と仮想、対抗項；LHAPDF との畳み込み & 固定次数の LO と NLO、重みの付け直し \\
$gg\to H$（有効理論） & 式~\eqref{eq:dgg}、\eqref{eq:dqgh} & 有効理論の1ループと実の放射、対抗項；畳み込み & HC\_NLO\_X0 と自作の NLO の UFO（R2 と紫外の対抗項を自分で導出）の MadLoop、木の振幅、固定次数 \\
MC@NLO & 4次元の実の放射（V$+$ジェット） & 自分の NLO と LO の V$+$ジェットとの比 & Pythia 8 のシャワー \\
$Z'$ の $K$ 因子と質量の下限 & 式~\eqref{eq:dqq}、\eqref{eq:dqg} & $Z'$ の結合での畳み込み & 自分の NLO の UFO の MadLoop と固定次数、段階 6 の上限 \\
\bottomrule
\end{tabularx}
\end{table}

\paragraph{一致の基準。}ヒッグスのノートと同じである。発散部分は $1/\epsilon$ の係数が記号として一致すること、有限部分は記号として一致すること、数値の積分は精度を明示する。\MGaMC{} の固定次数の積分は統計の誤差を持つので、全断面積とビンごとに、自分の値との差を誤差で割った値（プル）で比べる。$\mu_R$、$\mu_F$ の重みを付け直した値は同じ事象から求めるので、中心の値と統計の揺らぎが強く相関する。

\paragraph{手順。}次の順に進める。(1) Drell--Yan：$d$ 次元で式~\eqref{eq:real}、\eqref{eq:ff} を求めて極の打ち消しと式~\eqref{eq:dqq}、\eqref{eq:dqg} を確かめ、4次元の振幅の2乗を本ノートの FeynArts の模型と比べる。係数関数を LHAPDF のパートン分布と畳み込み、\MGaMC{} の固定次数の LO と NLO の全断面積、$m_{\ell\ell}$ の分布、9通りの $(\mu_R,\mu_F)$ と比べる（$Z$ の共鳴の付近、高い質量、$W^\pm$）。(2) $gg\to H$：重いトップの有効理論で、$gg$、$qg$、$q\bar q$ の係数関数を $d$ 次元で求める。NLO の UFO は、自分の模型から NLOCT で作ることを試み、使えなければ HC\_NLO\_X0 を使って、MadLoop の仮想の補正、実の放射の木の振幅、aMC@NLO の固定次数を確かめる。さらに、NLOCT が出せない頂点を自分で導いて自作の NLO の UFO を作り、同じ確かめをする。(3) MC@NLO：手順 (1)、(2) の過程で Pythia 8 のシャワーをつなぎ、全断面積が固定次数の NLO と一致すること、ラピディティの分布、横運動量の大きいところが LO の $V+$ジェットに近づくことを確かめる。(4) $Z'$：段階 6~\cite{env} の SSM の $Z'$ の NLO の $K$ 因子を、式~\eqref{eq:dqq}、\eqref{eq:dqg} の $Z'$ の結合での畳み込みと、ZpSSM の NLO の UFO の aMC@NLO で求め、質量の下限を NLO で出し直す。

\section{検算}\label{sec:checks}

手順 (1) の (a)〜(e) は、第~\ref{sec:coef} 節の係数関数を $d$ 次元で計算し直し、発散が消えること（第~\ref{sec:div} 節）と既知の形との一致を確かめる。(f) は係数関数とパートン分布から陽子の断面積を数値で求める部分の精度、(g) は \MGaMC{}（第~\ref{sec:mg5intro} 節）との比較、(h) は得られた断面積と $K$ 因子である。\MGaMC{} の値はモンテカルロの統計の誤差を持つので、差を誤差で割ったプルが $\pm2$ 程度に収まり、多くのビンの $\chi^2$ がビンの数と同じくらいなら一致とみなす。

手順 (1) は \texttt{Src/run\_dy\_check.sh} で再現できる（約3時間。ほとんどは \MGaMC{} の積分で、\texttt{--reuse} で既にある結果を使うと約30分。大半は自分のラピディティの分布の積分）。

\subsection{手順 (1)：Drell--Yan}\label{sec:step1}

\paragraph{(a) $d$ 次元の振幅の2乗。}\texttt{dy\_feyncalc.wl} は、$q\bar q\to V^*g$（2図）と $qg\to V^*q$（2図）の流れを FeynCalc で $d$ 次元で書き、スピンの和のトレースと両方のベクトルの偏極の和 $-g$ で $-g_{\mu\nu}H^{\mu\nu}$ を求める。Born と実の放射の両方で、軸性の流れ（$\gamma^\mu\gamma_5$、反可換の $\gamma_5$）とベクトルの流れの差は記号として0で、$Z$ と $W$ の $V-A$ の流れも同じ係数関数を持つ。Born に対する比は式~\eqref{eq:real} と記号として一致した（$-g_{\mu\nu}H_B^{\mu\nu}=4(1-\epsilon)\hat s$）。

\paragraph{(b) 位相空間の積分。}式~\eqref{eq:dr} の $v$ の積分は部分分数とベータ関数で $\epsilon$ について厳密に求め、$\epsilon=-0.3$、$z=0.37$ で \texttt{NIntegrate}（30桁）と15桁の範囲で一致した。$q\bar q$ の $(1-z)^2\times$（積分）は $z=1$ で正則で、$\delta(1-z)$、$\mathcal D_0$、$\mathcal D_1$ の係数と正則な部分に分けた（\ref{app:dy}）。

\paragraph{(c) 仮想の補正。}頂点の補正（Feynman ゲージ）を FeynCalc の TID で $B_0(Q^2)$ と $C_0(0,0,Q^2)$ に直し、閉じた形（時間的な $Q^2$ で $(-Q^2-i0)^{-\epsilon}=Q^{-2\epsilon}e^{i\pi\epsilon}$）を入れた。IBP の関係 $C_0=-\frac{2(d-3)}{(d-4)Q^2}B_0$ は閉じた形どうしで記号として成り立つ。軸性の流れとベクトルの流れの $F^{(1)}$ の差は0で、$F^{(1)}e^{-i\pi\epsilon}$ は実で、式~\eqref{eq:ff} の閉じた形と記号として一致した。

\paragraph{(d) 和と対抗項。}実の放射と仮想の補正と $\MSbar$ の共線の対抗項の和で、$q\bar q$ の $1/\epsilon^2$、$1/\epsilon$ と $qg$ の $1/\epsilon$ の係数はすべて記号として0になり、$\ln(\mu_R^2/Q^2)$ による微分も0だった。有限部分は式~\eqref{eq:dqq}、\eqref{eq:dqg}（Altarelli--Ellis--Martinelli の $\MSbar$ の形~\cite{aem79}、参照）と記号として一致した。ラピディティの分布の NLO（\ref{app:rap}）のために、$v$ を固定した部分 $R_0$、$R_1$、$S_0$、$S_1$ を同じ $d$ 次元の式から取り出し、$G$ が $v$ によらないときに式~\eqref{eq:dqq}、\eqref{eq:dqg} に記号として戻ることを確かめた。

\paragraph{(e) 4次元の振幅の2乗と模型。}NLO ノートの FeynArts の模型で、$u\bar u\to\gamma g$、$u\bar u\to Zg$、$u\bar d\to W^+g$、$ug\to\gamma u$（各2図、$V$ は $Q^2$ の off-shell）の4次元の $\sum|\mathcal M|^2$ を求めた。(a) の式の $\epsilon\to0$ に模型の結合（$eQ$、$\frac{g_Z}2(T_3-2Qs_W^2)$ と $\frac{g_Z}2T_3$、$\frac g{2\sqrt2}$ の $c_V$、$c_A$）と色の和 $N_cC_F$ を掛けたものと、4つとも記号として一致した。

\paragraph{(f) 自分の畳み込み。}\texttt{dy\_hadronic.py} は式~\eqref{eq:lumi} を数値で積分する。
\begin{itemize}
\item 係数関数を Python で書き直したものは、\texttt{dy\_feyncalc.wl} が7点の $z$ で出した値と $9\times10^{-16}$ で一致した。
\item $xf(x)$ は LHAPDF の値を $\ln x$ の4000点の3次スプラインにし、$G_{ij}$ は $\ln y$ の700点のスプラインにした。LHAPDF を直接呼んだ値との差は、$x<0.5$ の $xf$ で $8\times10^{-7}$、$G_{u\bar u}$ で $3\times10^{-6}$ である（$x\approx0.88$ では $b$ の $xf\approx10^{-9}$ で相対 $10^{-3}$ だが、効かない）。
\item $z$ の積分は $(\tau,\frac12)$ を $\ln z$ で、$(\frac12,1)$ を $1-z=u^3/2$ で Gauss--Legendre（96点と64点）にし、$Q^2$ の積分は共鳴の $\pm10\Gamma$ の中を Breit--Wigner の変数で（40点）、外を $\ln Q^2$ の区間（10~GeV より上は比1.5以下、下は比10以下、各20点）に分けて行った。点を倍にした変化は、66〜116~GeV の窓で $9\times10^{-8}$、全体の範囲（$Z$ は $m_{ee}>50$、500~GeV、$W$ は $10^{-3}$~GeV〜13~TeV）で $10^{-7}$ 以下だった。初めは $Q^2$ の全体を Breit--Wigner の変数だけで積分していて、$W$ の 50~GeV より下（全体の $1.3\%$）を少なく見積もり、全断面積が $10^{-3}$ 小さかった。また、\MGaMC{} は $m_{e\nu}$ の下限を持たないので、$W$ は $10^{-3}$~GeV から積分する（$m_{e\nu}<1$~GeV は $W^+$ の LO で 0.758~pb、$7.4\times10^{-5}$）。LHAPDF の $x_{\min}=10^{-9}$ より小さい $x$ は、\MGaMC{} と同じく LHAPDF の外挿（continuation）の値を使う。
\item ラピディティの分布は、$t=\ln x_1$ の400点の格子で $f(z,t)$ の累積の積分（Simpson）を作り、ビンの端で Hermite の補間（$C'=f$）をとって $G(z,v)$ を求めた（$v$ は48点）。$|y|<12$ のビンの和は、窓の断面積と LO で $-5\times10^{-7}$、NLO で $-1.2\times10^{-6}$ 一致し、$z$、$v$、$t$ の点を倍にしたビンの値の変化は $10^{-8}$ だった。
\end{itemize}

\paragraph{(g) \MGaMC{} の固定次数。}\texttt{dy\_mg5.sh} は、\texttt{loop\_sm-no\_b\_mass} で \texttt{p p > e+ e- [QCD]}、\texttt{p p > e+ ve [QCD]}、\texttt{p p > e- ve\textasciitilde{} [QCD]} を書き出し、固定次数（\texttt{fixed\_order = ON}）の LO と NLO を走らせる。$\mu_R=\mu_F=\mu_0$ は $Z$ の共鳴の付近（$m_{ee}>50$~GeV）で $M_Z$、高い質量（$m_{ee}>500$~GeV）で 1~TeV、$W$ で $M_W$ とし、$\mu_R$、$\mu_F$ を独立に $\frac12$、1、2 倍にした9通りを重みの付け直しで同じ事象から求めた。レプトンのカットはない（$Z$ は \texttt{mll\_sf} だけ）。分布は自分で書いた固定次数の解析 \texttt{analysis\_HwU\_dy.f}（$m_V$、$y_V$、$\log m_V$、質量の窓）で取った。積分の精度は \texttt{req\_acc\_FO}~$=3\times10^{-4}$ と、種を替えた $10^{-4}$（$W$ の NLO は $3\times10^{-4}$）の二組で、$W$ には裾の分布を加えた三組目（種を替えた）がある。表~\ref{tab:dytot} に中心のスケールの全断面積と窓を、表~\ref{tab:dychi} に分布の $\chi^2$ を示す。

\begin{table}[htbp]
\centering\small
\caption{自分の畳み込みと \MGaMC{} の固定次数（pb、$\mu_R=\mu_F=\mu_0$）。\MGaMC{} は独立な2〜3回の積分（1：$3\times10^{-4}$、2：種を替えた $10^{-4}$（$W$ の NLO は $3\times10^{-4}$）、3：$W$ の裾の分布の組）で、括弧は最後の桁の誤差、右の列はプル（\MGaMC{} から自分の値を引いて誤差で割ったもの）。$Z$ の全体は $m_{ee}>50$~GeV（高い質量は $m_{ee}>500$~GeV）、$W$ の全体はカットなし。}
\label{tab:dytot}
\begin{tabular}{llrrrrrrr}
\toprule
過程と範囲 & 次数 & 自分 & \multicolumn{2}{c}{\MGaMC{} 1} & \multicolumn{2}{c}{\MGaMC{} 2} & \multicolumn{2}{c}{\MGaMC{} 3} \\
\midrule
$Z$、全体 & LO & 1810.44 & 1810.16(42) & $-0.7$ & 1810.53(18) & $+0.5$ & &  \\
 & NLO & 2064.54 & 2063.80(67) & $-1.1$ & 2064.78(22) & $+1.1$ & &  \\
$Z$、66〜116~GeV & LO & 1721.13 & 1720.95(41) & $-0.4$ & 1721.04(17) & $-0.5$ & &  \\
 & NLO & 1966.68 & 1966.00(64) & $-1.1$ & 1966.83(21) & $+0.7$ & &  \\
$Z$、$m_{ee}>500$~GeV & LO & 0.0963761 & 0.0964092(190) & $+1.7$ & 0.0963768(74) & $+0.1$ & &  \\
 & NLO & 0.114451 & 0.114426(27) & $-0.9$ & 0.114448(10) & $-0.2$ & &  \\
$W^+$、全体 & LO & 10308.05 & 10307.7(22) & $-0.2$ & 10307.31(75) & $-1.0$ & 10307.89(75) & $-0.2$ \\
 & NLO & 11855.87 & 11863.0(22) & $+3.3$ & 11856.9(20) & $+0.5$ & 11850.5(31) & $-1.7$ \\
$W^+$、60〜100~GeV & LO & 10005.38 & 10005.2(21) & $-0.1$ & 10004.51(73) & $-1.2$ & 10005.02(73) & $-0.5$ \\
 & NLO & 11543.41 & 11545.0(16) & $+1.0$ & 11543.1(12) & $-0.3$ & 11542.7(13) & $-0.6$ \\
$W^-$、全体 & LO & 7776.65 & 7776.7(17) & $+0.0$ & 7776.44(57) & $-0.4$ & 7776.47(57) & $-0.3$ \\
 & NLO & 8853.23 & 8856.0(31) & $+0.9$ & 8856.5(19) & $+1.8$ & 8853.8(20) & $+0.3$ \\
$W^-$、60〜100~GeV & LO & 7541.22 & 7541.1(16) & $-0.1$ & 7541.14(55) & $-0.1$ & 7540.96(55) & $-0.5$ \\
 & NLO & 8616.52 & 8616.3(17) & $-0.1$ & 8618.14(95) & $+1.7$ & 8616.26(90) & $-0.3$ \\
\bottomrule
\end{tabular}
\end{table}

\begin{table}[htbp]
\centering\small
\caption{分布の $\chi^2$（中心のスケール、独立な積分の和）と自由度。$y$ の分布は $Z$ で 66〜116~GeV（高い質量は $m_{ee}>500$~GeV で、$|y|\ge3.25$ の運動学の端の断面積 $10^{-13}$~pb のビンを除く）、$W$ で全体。$\log_{10}m_V$ は $W$ の三組目だけ（$10^{-3}$~GeV〜13~TeV の0.2刻み）。}
\label{tab:dychi}
\begin{tabular}{lcccccc}
\toprule
 & \multicolumn{2}{c}{$m_V$ のビン} & \multicolumn{2}{c}{$y_V$ のビン} & \multicolumn{2}{c}{$\log_{10}m_V$} \\
過程 & LO & NLO & LO & NLO & LO & NLO \\
\midrule
$Z$（50〜150~GeV、5~GeV 刻み） & 36.9/40 & 30.5/40 & 74.2/80 & 74.0/80 & — & — \\
$Z$（0.5〜3.5~TeV、200~GeV 刻み） & 28.0/30 & 32.8/30 & 41.1/52 & 52.7/52 & — & — \\
$W^+$（50〜150~GeV） & 44.3/60 & 47.3/60 & 109.8/120 & 131.1/120 & 19.8/27 & 30.0/27 \\
$W^-$（50〜150~GeV） & 62.0/60 & 49.3/60 & 104.0/120 & 128.5/120 & 17.6/28 & 25.0/28 \\
\bottomrule
\end{tabular}
\end{table}

$Z$ と $W$ の窓、$m_V$ の分布、ラピディティの分布は、LO と NLO で統計の範囲で一致した。9通りの $(\mu_R,\mu_F)$ の重みを付け直した値も、$W$ の NLO の全体を除くすべての全断面積で自分の値と $1.9\sigma$ 以内で合った（$\mu_R$ は NLO の補正の $\alpha_s(\mu_R)$ だけを、$\mu_F$ はパートン分布と $L_F$ の項を変える）。$W$ の NLO の全体だけ、3回の積分のばらつきが誤差より大きい（$W^+$ で平均のまわりの $\chi^2=11.5$、自由度2。9通りのスケールでは最大 $4.9\sigma$）。裾の分布で、そのばらつきは $m_{e\nu}<1$~GeV にあることが分かった。そこでは $\mu_F=M_W$ の $L_F=\ln(\mu_F^2/Q^2)\gtrsim9$ の NLO の補正が LO（$W^+$ で 0.758~pb）より大きく負で（自分の NLO は $-0.629$~pb）、正と負の重みが打ち消し合う。\MGaMC{} の値は $W^+$ で $-4.93\pm2.5$~pb、$W^-$ で $-0.14\pm1.3$~pb（自分は $-0.620$~pb）で、この 1~GeV より下の積分が全体の誤差（2〜3~pb）のほとんどを占める。1~GeV より上の裾（1〜10、10〜50、150~GeV〜13~TeV）は LO と NLO で $2.1\sigma$ 以内で一致した。

\paragraph{(h) 断面積と $K$ 因子。}表~\ref{tab:dyk} に自分の値をまとめる。共鳴の付近の $K$ は 1.14〜1.15 で、$q\bar q$ のチャンネルが $+29$〜$30\%$、$qg$ が $-14$〜$-16\%$ である。高い質量では $qg$ の光度が相対的に小さく $K=1.19$ になる。$\mu_R$、$\mu_F$ の9通りの範囲は、LO（$\mu_F$ だけ）の $\pm12$〜$13\%$ から NLO で $^{+4}_{-7}\%$ に縮む。

\begin{table}[htbp]
\centering\small
\caption{自分の LO と NLO の断面積（pb、$\mu_0$ は (g) の値）、$K$、NLO の補正のチャンネルごとの分（LO に対する比）、9通りのスケールの範囲。}
\label{tab:dyk}
\begin{tabular}{lrrcccc}
\toprule
過程 & LO & NLO & $K$ & $q\bar q$ & $qg$ & スケールの範囲（NLO；LO） \\
\midrule
$Z$、$m_{ee}>50$~GeV & 1810.44 & 2064.54 & 1.1404 & $+0.2940$ & $-0.1537$ & $^{+3.6}_{-6.4}\%$；$^{+11.4}_{-12.4}\%$ \\
$Z$、$m_{ee}>500$~GeV & 0.0963761 & 0.114451 & 1.1875 & $+0.2838$ & $-0.0962$ & $^{+1.6}_{-1.2}\%$；$^{+2.1}_{-2.2}\%$ \\
$W^+$ & 10308.05 & 11855.87 & 1.1502 & $+0.2945$ & $-0.1444$ & $^{+4.0}_{-7.0}\%$；$^{+12.0}_{-13.0}\%$ \\
$W^-$ & 7776.65 & 8853.23 & 1.1384 & $+0.3004$ & $-0.1620$ & $^{+3.7}_{-6.9}\%$；$^{+12.3}_{-13.4}\%$ \\
\bottomrule
\end{tabular}
\end{table}

測定値（参照）と並べると、ATLAS の 13~TeV の全断面積~\cite{atlaswz}（1レプトンのフレーバーあたり）$\sigma_Z(66<m_{\ell\ell}<116~\mathrm{GeV})=1.981\pm0.057$~nb、$\sigma_{W^+}=11.83\pm0.40$~nb、$\sigma_{W^-}=8.79\pm0.30$~nb（誤差は統計、系統、ルミノシティの2乗和）に対し、本ノートの NLO は 1.967、11.856、8.853~nb で、比は 0.993、1.002、1.007 である。ただし本ノートの $W$ は CKM 行列を単位行列とした値である。PDG の $|V_{ij}|$（参照）を入れると（\texttt{dy\_ckm.py}。$u\bar s$、$c\bar d$ などの組が加わり、$u\bar d$、$c\bar s$ は $|V_{ud}|^2$、$|V_{cs}|^2$ に減る。各行の $|V_{ij}|^2$ の和は1だが、入れ替わる反クォークのパートン分布が違う）、$W^+$、$W^-$ の NLO は $0.35\%$、$0.62\%$ 下がって 11.814、8.798~nb になり、比は 0.999、1.001 になる（LO でも $0.33\%$、$0.57\%$ で、QCD の補正はフレーバーによらない）。$Z$ は中性の流れなので CKM によらない。\MGaMC{} との比較は、両方が単位行列を使うので、この違いを含まない。本ノートは NNLO の QCD と電弱の補正を入れていないので、この一致は NLO の範囲のもので、$\mu$ の範囲（約 $5\%$）より小さい。

\paragraph{(i) 道具の注意。}\MGaMC{} の固定次数の解析の HwU の値はビンごとの断面積で、幅で割っていない。NLO の run\_card の固定スケールの既定値は 91.118~GeV（$M_Z$ ではない）。\texttt{pdlabel = lhapdf} では $\alpha_s(Q)$ も LHAPDF の値で（\texttt{alfas\_functions\_lhapdf6.f}）、自分の計算も同じ関数を使った。$W$ の NLO は $m_{e\nu}$ の下限がなく、$3\times10^{-4}$ の精度に約30分かかる（$Z$ は1.5分）。重みを付け直した列の誤差は HwU にないので、中心の列の誤差に $\max(1,|\text{付け直した値}/\text{中心の値}|)$ を掛けた（値の比だけを掛けると、NLO が LO を打ち消す $W$ の数 GeV のビンで誤差が小さくなりすぎる）。これらは計算環境の記録~\cite{env} の表に加えた。
\subsection{手順 (2)：重いトップの有効理論の $gg\to H$}\label{sec:step2}

手順 (2) は \texttt{Src/run\_gh\_check.sh} で再現できる（約1時間10分。ほとんどは \MGaMC{} の積分で、\texttt{--reuse} で既にある結果を使うと約2分。NLOCT の対抗項の計算 (d) は \texttt{--nloct} を付けたときだけ走らせ（約15分）、ふだんは保存したその結果から自作の NLO の UFO (i) を作る。HC\_NLO\_X0 は \nolinkurl{~/hep/src/smpp2-dl/extracted} に要る）。この手順の \MGaMC{} の模型は、FeynRules の模型のデータベースの HC\_NLO\_X0（版 1.3~\cite{hcnlo}、\texttt{HC\_NLO\_X0\_UFO.zip}、123{,}075 バイト、\nolinkurl{~/hep/src/smpp2-dl}）で、取得は承認を得て行った。その $ggX0$ の結合は $-i\cos\alpha\,\kappa_{Hgg}g_{Hgg}$、$g_{Hgg}=-g_s^2/(12\pi^2v)=-\alpha_s/(3\pi v)$ で、$\kappa_{Hgg}=\cos\alpha=1$ で本ノートの $c_{Hgg}/v$ に等しい。

\paragraph{(a) 振幅の2乗。}\texttt{gh\_feyncalc.wl} は、ヒッグスのノートの FeynArts の模型（\texttt{SM\_CKM\_HEFT\_FA}）で、$gg\to H$ の Born（1図）と、$gg\to Hg$（4図）、$qg\to Hq$（3図）、$q\bar q\to Hg$（3図）の $d$ 次元の振幅の2乗を、グルーオンの物理的な偏極の和（基準のベクトルは別の運動量）で求める。4次元では、それぞれ $(M_H^8+\hat s^4+t^4+u^4)/(\hat stu)$、$(\hat s^2+u^2)/t$、$(t^2+u^2)/\hat s$ の定数倍になった。

\paragraph{(b) 位相空間と仮想の補正。}Drell--Yan と同じ $v$ の積分と分布の展開で実の放射を求め、仮想の補正（$c_{Hgg}g_s^2$ の次数の14図、グルーオンとゴーストのループ）は TID で $B_0(M_H^2)$ と $C_0$ にし、閉じた形を入れた。$F^{(1)}e^{-i\pi\epsilon}$ は実だった。

\paragraph{(c) 和と対抗項。}紫外と共線の対抗項を足すと、$gg$ の2重極と単極、$qg$ と $q\bar q$ の単極はすべて記号として0になり、$L_R$ による微分は $\Delta_{gg}$ の $\delta(1-z)$ の $2\beta_0$ だけだった（LO の $\alpha_s^2$ の走りを打ち消す）。$z\Delta^{(1)}$ は三つのチャンネルとも文献の形~\cite{dawson91,dsz91,hk02}（参照）と記号として一致した。

\paragraph{(d) NLOCT による NLO の UFO（そのままでは使えなかったこと）。}NLO の UFO を自分で作るため、QCD ノートの模型から光子、電子、重いクォークを除き、ヒッグスと演算子を足した \texttt{QCDH.fr}（演算子の係数を $\alpha_s$ で書き、$\alpha_s$ の繰り込みが演算子の対抗項を与えるようにした）を書き、QCD ノートと同じ NLOCT の三つの手順（\texttt{gh\_generate\_nlo.sh}）を通した。FeynArts の \texttt{Lorentz.gen} にはスカラー1個とベクトル3、4個の結合の形がないので、FeynRules の generic のファイルを NLOCT に渡すと計算は終わったが、$GGH$ と $q\bar qH$ の R2 と UV の頂点に Feynman のパラメータの積分と内部の運動量が残り、使えなかった（NLOCT 1.02 は次元5の演算子の頂点を扱わない。計算環境の記録~\cite{env} の表）。そこでまず HC\_NLO\_X0 を使い、その R2 と UV の頂点を (e) で確かめる。のちに、この二つの頂点を自分で導いて NLOCT の $H$ の頂点と入れ替え、自作の NLO の UFO を作った（(i)）。

\paragraph{(e) MadLoop の仮想の補正。}\texttt{gh\_madloop.py} は \texttt{g g > x0 / t [virt=QCD]} を standalone で書き出し、$M_H=125.2$~GeV、$\alpha_s=0.118$、$\mu_R=M_H/2$、$M_H$、$2M_H$ で、MadLoop の規格化（$(4\pi)^\epsilon/\Gamma(1-\epsilon)$ を外し、$\ln(\mu_R^2/M_H^2)$ を係数に含む）の $A/\epsilon^2+B/\epsilon+C$ を自分の値（\texttt{gh\_feyncalc.wl} の (f)：$\MSbar$ の $(4\pi)^\epsilon e^{-\epsilon\gamma_E}$ との比 $\Gamma(1-\epsilon)e^{-\epsilon\gamma_E}$ を掛ける）と比べた。$\mu_R=M_H$ で自分の値は $A=-12$、$B=-2\beta_0=-15.333$、$C=6\pi^2=59.218$ である。初め、$A$ は $10^{-14}$ で一致したが、$B$ は $\mu_R$ によらず $-4.000$、$C$ は $+22.000$ ずれた。$C$ のずれは、HC\_NLO\_X0 の $ggX0$ の UV の対抗項の有限部分に入っているウィルソン係数の補正 $2\times\frac{11}4\frac{\alpha_s}\pi=22a_s$ である。$B$ のずれは、その対抗項の軽いクォーク（$u$、$d$、$s$）の部分 $3\times\frac43a_s$ が、ループの粒子として \texttt{[[u],[d],[s]]} に結び付けてあるため、軽いクォークのループのない $gg\to X0$ では MG5 3.8.1 に落とされることによる。作業用の複製でこの極を \texttt{[[G]]} の部分に移すと（有限部分は変わらない）、$A$、$B$、$C+22$ は三つの $\mu_R$ で $2\times10^{-14}$ 以内で一致した。これで HC\_NLO\_X0 の R2 の頂点と UV の対抗項が確かめられた。直す前は、aMC@NLO が \texttt{Poles do not cancel} で止まった（FKS の単極 $-38.23$ に対し MadLoop $-42.23$、2重極の $\frac1{12}$ を単位として）。

\paragraph{(f) 実の放射の木の振幅。}\texttt{gh\_tree.py} は \texttt{g g > x0 g}、\texttt{u g > x0 u}、\texttt{u u\textasciitilde{} > x0 g} の MadGraph の振幅の2乗を、(a) の4次元の式と、それぞれ3点で比べ、$7\times10^{-16}$ で一致した。

\paragraph{(g) \MGaMC{} の固定次数。}\texttt{gh\_mg5.sh} は \texttt{p p > x0 / t [QCD]}（\texttt{HC\_NLO\_X0-heft}、5フレーバー、上の直しを入れた作業用の複製）の LO と NLO を、$M_H=125.2$~GeV、NNPDF4.0 NLO、$\mu_R=\mu_F=M_H$ で走らせる（\texttt{req\_acc\_FO}~$=3\times10^{-4}$ と、種を替えた $10^{-4}$）。比べるときに分かったことが二つある。一つは、HC\_NLO\_X0 の $ggX0$ の結合の次数が \texttt{QNP} だけで、その $\alpha_s^2$ が $g_{Hgg}$ の中にあるため、\MGaMC{} の $\mu_R$ の重みの付け直しが $\alpha_s^2$ を動かさないことである（LO が $\mu_R$ によらなかった）。$\mu_F$ は重みの付け直しで、$\mu_R$ は中心のスケールを $M_H/2$、$2M_H$ にした別の実行で比べた。もう一つは、FKS が実の放射の過程を Born（\texttt{g g > x0}）の脚を分けて作るので、赤外で有限な $q\bar q\to Hg$ が \texttt{[QCD]} の過程に入らないことである（実の放射の過程の一覧に $gg\to Hg$ と $qg$、$gq$ だけがあった）。初め、NLO が自分の値より $5\times10^{-4}$ 小さく（$-3.9\sigma$、$-6.0\sigma$）、その差 $0.0153\pm0.0022$~pb が自分の $q\bar q$ のチャンネル 0.0174~pb と合ったことで分かった。$q\bar q\to Hg$ は MadEvent の LO で別に求め、自分の $q\bar q$ の部分と比べた。結果を表~\ref{tab:ghtot} に示す。すべて統計の範囲で一致した。

\begin{table}[htbp]
\centering\small
\caption{$pp\to H$（重いトップの有効理論、ウィルソン係数の補正を含む）の自分の値と \MGaMC{}（pb）。NLO の自分の値は $q\bar q$ のチャンネルを除いたもの（\MGaMC{} の \texttt{[QCD]} に入らない）。括弧はプル。}
\label{tab:ghtot}
\begin{tabular}{lllll}
\toprule
$(\mu_R,\mu_F)/M_H$ & 次数 & 自分 & \MGaMC{} 1 & \MGaMC{} 2 \\
\midrule
$(1,1)$ & LO & 12.9244 & 12.9238(26) ($-0.2$) & 12.9247(9) ($+0.4$) \\
$(1,2)$ & LO & 13.4039 & 13.4033(27) ($-0.2$) & 13.4043(9) ($+0.4$) \\
$(1,\frac12)$ & LO & 12.2229 & 12.2223(26) ($-0.2$) & 12.2232(9) ($+0.3$) \\
$(\frac12,1)$、$(2,1)$ & LO & 15.9485、10.6929 & 15.9471(33) ($-0.4$)、10.6919(22) ($-0.4$) & \\
$(1,1)$ & NLO & 29.5852 & 29.5777(64) ($-1.2$) & 29.5886(23) ($+1.5$) \\
$(1,2)$ & NLO & 30.2182 & 30.2099(65) ($-1.3$) & 30.2226(24) ($+1.9$) \\
$(1,\frac12)$ & NLO & 28.7947 & 28.7880(64) ($-1.1$) & 28.7973(23) ($+1.1$) \\
$(\frac12,1)$ & NLO & 35.4113 & 35.4270(83) ($+1.9$) & \\
$(2,1)$ & NLO & 25.0834 & 25.0857(54) ($+0.4$) & \\
\midrule
$q\bar q\to Hg$（LO） & & 0.017500 & 0.017471(20) ($-1.5$) & \\
\bottomrule
\end{tabular}
\end{table}

LO のヒッグスのラピディティの分布（38 ビン）は $\chi^2=15.7$、20.0 で一致した。$\mu_R$ を替えた実行のうち $\mu_F$ を重みで動かしたものも、最大 $2.4\sigma$ で合った。これで、$gg\to H$ の \MGaMC{} の NLO は、仮想の補正（R2 と UV を含む）、実の放射、FKS の引き算、パートン分布との畳み込み、$\mu_F$ の重みの付け直しのすべてで確かめられた。ただし (e) と上の二つの点に注意が要る。

\paragraph{(h) 断面積と $K$ 因子。}$\mu=M_H$ で、LO は 12.924~pb、NLO は 29.603~pb（ウィルソン係数の補正を含む。含まないと 27.054~pb）で、$K=2.29$ である。NLO の補正の内訳は、LO に対して $gg$ のチャンネルが $+111\%$、ウィルソン係数が $+19.7\%$、$qg$ が $-1.8\%$、$q\bar q$ が $+0.13\%$ で、ほとんどがグルーオンの放射である。9通りのスケールの範囲は NLO で $^{+22}_{-18}\%$ で、LO（$\mu_R=\mu_F$）の $^{+17}_{-14}\%$ より狭くならない。$gg\to H$ では NLO の補正が LO と同じくらい大きく、さらに高い次数（NNLO、N$^3$LO）が必要なことを示している。トップの質量を正しく入れた LO との比 1.066 を掛けると 31.56~pb になる。LHC Higgs Cross Section Working Group の N$^3$LO の値~\cite{yr4ggh}（参照、$M_H=125$~GeV、13~TeV で 48.58~pb）とは、本ノートが入れていない NNLO、N$^3$LO の QCD の補正と電弱の補正、スケールの選び方（参照は $\mu=M_H/2$）の分だけ違う。


\paragraph{(i) 自作の NLO の UFO。}MadLoop は、ループの積分の被積分関数の分子を、ループの運動量 $q$ の4次元の多項式として数値で組み立て、既知のスカラーの積分の和に分ける（OPP の方法）。ところが次元正則化では、ループの運動量は $D=4-2\epsilon$ 次元にある。$q$ を4次元に限ると、分子の中の $\bar q^2=q^2+\tilde q^2$（$\tilde q$ は余分な $-2\epsilon$ 次元の成分）の $\tilde q^2$ と、計量の縮約 $g^\mu{}_\mu=D$ の $D-4$ が落ちる。これらは $\epsilon$ に比例するが、ループの紫外の発散 $1/\epsilon$ に掛かると有限な項が残る。この残りを R2 と呼び、MadLoop は UFO に書かれた木の形の頂点（R2 の頂点）として足す。紫外の対抗項も UFO の頂点である。NLOCT はこの二つを繰り込み可能な頂点についてしか出せない（(d)）ので、$GGH$ の頂点のものを自分で導いた。

\texttt{gh\_r2.wl} は、ヒッグスのノートの FeynArts の模型で、1ループの $G^a_\mu(k_1)G^b_\nu(k_2)\to H$ の頂点（$c_{Hgg}g_s^2$ の次数の14図、Feynman ゲージ。外の添字と運動量は4次元、ループは $D$ 次元の 't~Hooft--Veltman の方式）を、分子を $D$ 次元で計算したものから4次元で計算したものを引いて求める。グルーオンは質量殻から外す（$k_1^2,k_2^2\neq0$）ので赤外の発散はなく、差は紫外の極だけから来る。$\tilde q^2$ を含む積分は、$\int\tilde q^2Y=(D-4)\times$（$\int q^\alpha q^\beta Y$ の $g^{\alpha\beta}$ の係数）、$\tilde q^4$ は $(D-4)(D-2)\times$（$g^{\alpha\beta}g^{\gamma\delta}$ の係数）で求めた（$\tilde q$ は4次元の外の運動量と直交する）。結果は
\begin{equation}
R_2^{\mu\nu,ab}=\frac{ic_{Hgg}}v\frac{g_s^2}{128\pi^2}\delta^{ab}\bigl[-17k_1^2g^{\mu\nu}-93\,k_1\!\cdot\!k_2\,g^{\mu\nu}-17k_2^2g^{\mu\nu}+14k_1^\mu k_1^\nu+k_1^\mu k_2^\nu+89k_2^\mu k_1^\nu+14k_2^\mu k_2^\nu\bigr]
\label{eq:r2ggh}
\end{equation}
で、すべて $\tilde q^2$ の項から来た（$D-4$ の項の和は0だった）。これは HC\_NLO\_X0 の R2 の頂点（参照）と、7つの係数がすべて一致した。紫外の対抗項は、演算子の係数が $\alpha_s$ に比例するので、$\alpha_s$ の $\MSbar$ の繰り込み $\delta\alpha_s/\alpha_s=-\beta_0\frac{g_s^2}{16\pi^2}\frac1\epsilon$（$\beta_0=11-\frac23n_f=\frac{23}3$）を木の頂点 $i(c_{Hgg}/v)\delta^{ab}(k_1\!\cdot\!k_2\,g^{\mu\nu}-k_2^\mu k_1^\nu)$ に掛けたものである（質量0のグルーオンの場の繰り込みは質量殻の方式で0）。これは (e) で直した HC\_NLO\_X0 のものと一致した。木の頂点の形も、FeynArts から求めたものが手で書いた形と一致した。

\texttt{gh\_nlo\_ufo.wl} は、NLOCT の対抗項から $H$ を含む頂点（$GGH$ は積分が残り、$q\bar qH$ は内部の運動量が残り、$GGHH$ は $\kappa_{Hgg}^2$ の次数）を除き、この二つを入れて、FeynRules の \texttt{WriteUFO} で \texttt{generated/QCDH\_NLO\_UFO} を書く（$c_{Hgg}=-\frac{\alpha_s}{3\pi}\kappa_{Hgg}$）。二つの頂点のループの粒子はグルーオンだけにした。$gg\to H$ の1ループの図にはグルーオンのループしかないので、$n_f=5$ の $\delta\alpha_s$ の全体をグルーオンに結び付ければ、(e) のように MG5 に落とされない。ウィルソン係数の補正は入れていない（頂点は演算子そのもの）。FeynRules が書いた木の頂点は FeynCalc のものと一致した。$q\bar qH$ の頂点は、木がなく、$gg\to H$ の1ループに現れないので要らない。

\texttt{gh\_own\_madloop.py} の MadLoop（\texttt{g g > h [virt=QCD]}）は、何も直さずに、Born が $(c_{Hgg}/v)^2M_H^4/64$ と $4\times10^{-16}$ で、$A$、$B$、$C$ が自分の値と三つの $\mu_R$ で $2\times10^{-14}$ 以内で一致した（ウィルソン係数を入れていないので、$C+22$ ではなく $C$ と一致する）。これで、自分で導いた R2 と紫外の対抗項が MadLoop の中で正しく働くことが確かめられた。

\texttt{gh\_own\_mg5.sh} の aMC@NLO（\texttt{p p > h QCD=2 HIG=1 [QCD]}、$\kappa_{Hgg}=1$、\texttt{req\_acc\_FO}~$=10^{-4}$）を動かすには、ツールの側の手当てが二つ要った。結合の次数を書かないと、aMC@NLO は模型の次数の階層から QED と QCD の次数を推し量ろうとして \texttt{KeyError: 'QED'} で止まる（この模型には QED の次数がない）ので、次数を書いた。また、FKS の \texttt{fks\_singular.f} が $M_Z$、$M_W$ と、使わない6→5フレーバーの対抗項の $m_t$ を参照してコンパイルが止まるので、作業用の複製に頂点のない $Z$、$W$、$t$ と、どの頂点も持たない QED の次数を足した（QCD ノートと同じ。計算環境の記録~\cite{env} の表）。結果を表~\ref{tab:ghown} に示す。自分の値はウィルソン係数の補正と $q\bar q$ のチャンネルを除いたもので、$v$ は模型の値 $246.21965$~GeV（HC\_NLO\_X0 は $246.21846$~GeV）である。今度は $\alpha_s^2$ が結合の次数 \texttt{QCD=2} として入っているので、$\mu_R$ の重みの付け直しも効く。9通りの $(\mu_R,\mu_F)$ のすべてで、LO は $+1.1\sigma$〜$+1.4\sigma$、NLO は $-1.3\sigma$〜$+0.1\sigma$ で一致した（同じ事象の重みを付け直した9つの値は互いに相関している）。LO のヒッグスのラピディティの分布（38 ビン）は $\chi^2=25.0$（最大のプル 1.9）だった。aMC@NLO の初めの極の確かめも、直しなしで通った（20点のすべて）。これで、配布の模型に頼らず、自分の模型と自分で導いた R2 と紫外の対抗項から作った NLO の UFO で、$gg\to H$ の \MGaMC{} の NLO（仮想の補正、実の放射、FKS の引き算、9通りのスケールの重みの付け直し）が確かめられた。

\begin{table}[htbp]
\centering\small
\caption{$pp\to H$（重いトップの有効理論、ウィルソン係数の補正なし）の自分の値と、自作の NLO の UFO の \MGaMC{}（pb、\texttt{req\_acc\_FO}~$=10^{-4}$、9通りは重みの付け直し）。NLO の自分の値は $q\bar q$ のチャンネルを除いたもの。括弧はプル。}
\label{tab:ghown}
\begin{tabular}{lllll}
\toprule
$(\mu_R,\mu_F)/M_H$ & LO 自分 & LO \MGaMC{} & NLO 自分 & NLO \MGaMC{} \\
\midrule
$(1,1)$ & 12.9243 & 12.9256(9) ($+1.4$) & 27.0366 & 27.0353(23) ($-0.5$) \\
$(1,2)$ & 13.4038 & 13.4051(9) ($+1.4$) & 27.5750 & 27.5723(24) ($-1.1$) \\
$(1,\frac12)$ & 12.2228 & 12.2240(9) ($+1.4$) & 26.3844 & 26.3846(23) ($+0.1$) \\
$(\frac12,1)$ & 15.9483 & 15.9499(11) ($+1.4$) & 31.9178 & 31.9161(27) ($-0.6$) \\
$(\frac12,2)$ & 16.5400 & 16.5416(11) ($+1.4$) & 32.4651 & 32.4614(28) ($-1.3$) \\
$(\frac12,\frac12)$ & 15.0826 & 15.0841(10) ($+1.4$) & 31.3030 & 31.3033(27) ($+0.1$) \\
$(2,1)$ & 10.6927 & 10.6938(9) ($+1.2$) & 23.1655 & 23.1646(23) ($-0.4$) \\
$(2,2)$ & 11.0895 & 11.0905(9) ($+1.2$) & 23.6753 & 23.6733(23) ($-0.9$) \\
$(2,\frac12)$ & 10.1124 & 10.1134(9) ($+1.1$) & 22.5217 & 22.5219(23) ($+0.1$) \\
\bottomrule
\end{tabular}
\end{table}

\subsection{手順 (3)：MC@NLO と Pythia 8}\label{sec:step3}

\paragraph{考え方。}固定次数の NLO は、グルーオンを高々1個放つところまでを正しく数える。実際の衝突では、始状態のパートンは陽子から出てくるまでに何度もグルーオンを放ち、その反跳でボソンは横向きの運動量 $p_T$ を受け取る。$p_T$ が小さいと、1回の放射の確率 $\sim\alpha_s\ln^2(Q^2/p_T^2)$ が 1 に近づくので、何回もの放射を足し上げなければならない。\textbf{パートンシャワー}（ここでは Pythia~8）は、「ある $p_T$ より上では何も放たない確率」（\textbf{Sudakov の因子}）を使って放射を順に確率的に作り、この足し上げを近似する。固定次数の NLO の $p_T$ の分布は $p_T\to0$ で発散する（打ち消しの相手の仮想の補正が $p_T=0$ にある）が、シャワーをつなぐと有限な分布になる。

NLO の事象にそのままシャワーをかけると、1個目の放射を NLO の実の放射とシャワーの両方で数えてしまう。\textbf{MC@NLO}~\cite{mcatnlo} は、シャワーが1回目に放つ分（MC の対抗項。シャワーの放射の確率を $\alpha_s$ の1次で展開したもの）を NLO の実の放射から引き、同じものを Born の事象の側に足しておく。こうすると、シャワーをかけた結果を $\alpha_s$ の1次で展開したものが NLO と一致し、それより高い次数をシャワーが補う。引き算のため、事象の重みには負のものが混じる。シャワーは確率を保つ（事象を消さず、放射を足すだけ）ので、次のことが期待される。
\begin{itemize}
\item[(i)] 全断面積は固定次数の NLO と統計の範囲で一致する。
\item[(ii)] ボソンのラピディティのように放射に鈍い量は、NLO と $\alpha_s^2$ の違いしかない。
\item[(iii)] $p_T$ が大きいところは1回の硬い放射で決まり、LO の V$+$ジェット（$\alpha_s$ の1次の実の放射）に近づく。
\item[(iv)] $p_T$ が小さいところは、発散のかわりに Sudakov の山になる。
\end{itemize}
手順 (3) ではこれらを自分の計算と比べる。

手順 (3) は \texttt{Src/run\_mc\_check.sh} で再現できる（約5分。手順 (1)、(2) の固定次数の実行が要る）。

\paragraph{(a) 自分の LO の V$+$ジェットと H$+$ジェット。}比べる相手の大きい $p_T$ の分布は、手順 (1)、(2) で求めた4次元の実の放射から作る（\texttt{mc\_hadronic.py}）。Drell--Yan では、$t=-\hat s(1-z)v$、$u=-\hat s(1-z)(1-v)$ として
\[
\frac{d\sigma_{ij}}{dQ^2\,dv}=\frac{\sigma_0(Q^2)}{\hat s}\,a_s(1-z)\,r_{ij}(z,v),\qquad r_{q\bar q}=2C_F\Bigl(\frac ut+\frac tu+\frac{2Q^2\hat s}{ut}\Bigr),\quad r_{qg}=-2T_R\Bigl(\frac u{\hat s}+\frac{\hat s}u+\frac{2Q^2t}{u\hat s}\Bigr)
\]
で、ボソンの横運動量は $p_T^2=tu/\hat s=\hat s(1-z)^2v(1-v)$ である。$p_T$ のビンに入る $v$ の範囲で積分し、Drell--Yan の Born と同じ $Q^2$ と $z$ の畳み込みをする（$p_T>0$ ではプラスの分布は要らない）。$gg\to H$ では \texttt{gh\_tree.py} の $gg\to Hg$、$qg\to Hq$、$q\bar q\to Hg$ の振幅の2乗を使う。積分の節点を倍にしても $10^{-7}$ しか動かない。固定次数の NLO の $p_T>0$ のビンには実の放射の事象しか入らない（Born、仮想の補正、引き算の事象は $p_T=0$）ので、手順 (1)、(2) の固定次数の実行の $p_T$ の分布（5~GeV ごとに 250~GeV まで）でこの式を確かめられる。$p_T>5$~GeV の49ビンで、$Z$ は $\chi^2=59.7$、43.0（2回の積分）、和は自分の 1485.35~pb に対し 1485.26~pb、1484.88~pb、$H$（$q\bar q$ を除く）は $\chi^2=37.3$、自分の 26.993~pb に対し 26.994~pb で一致した。

\paragraph{(b) MC@NLO の実行。}\texttt{mc\_mg5.sh} は、$Z$ を \texttt{loop\_sm-no\_b\_mass} の \texttt{p p > e+ e- [QCD]}（$m_{ee}>50$~GeV、$\mu_R=\mu_F=M_Z$）で、$H$ を自作の NLO の UFO の \texttt{p p > h [QCD]}（$\mu_R=\mu_F=M_H$、ウィルソン係数の補正なし）で、それぞれ20万事象作り、Pythia~8.317~\cite{pythia83} でシャワーをかける。比べやすさのため、パートンの段階で止める（ハドロン化、多重パートン相互作用、QED のシャワー、primordial $k_T$ を切る）。解析は自分で書いた \texttt{py8an\_HwU\_mc.f} で、固定次数と同じヒストグラム（ボソンのラピディティと $p_T$）を作る。$Z$ のボソンは $e^+e^-$ の対である（始状態の放射の反跳は対を丸ごと動かすので、$m_{ee}$ は変わらない）。

動かすまでに、ツールの側の問題が三つあった（計算環境の記録~\cite{env} の表）。
\begin{itemize}
\item シャワーのプログラムのリンクが \texttt{library 'gfortran' not found} で止まった。shower\_card の \texttt{EXTRAPATHS} に libgfortran の場所を足して直した。
\item スケールの重みを切ると、事象に重みの塊（\texttt{<rwgt>}）がなく、\MGaMC{} の事象の読み込み（\texttt{LHEFRead.h}）が終わらないループに入った。そのループは1回ごとに LHE のヘッダーを付け足すので、シャワーの1つの仕事が数十 GB のメモリを使い、10並列でマシンが止まった。スケールの重みを入れて直した（9通りのスケールの重みは、シャワーのあとの分布にも付く）。
\item 自作の模型では、QCD のほかの結合の次数が \texttt{HIG} だけなので、事象を作る段の \texttt{split\_type} が「QED の位置 $=-1$」の配列の外を読み、\texttt{IPROCMAP: NOT IMPLEMENTED} で止まった。MC@NLO 用の複製で、次数の名前を \texttt{QED} に替えて直した。
\end{itemize}
直したあとは、シャワーの1つの仕事のメモリは約 0.4~GB、全体で約5分だった。


\begin{figure}[htbp]
\centering
\begin{tikzpicture}
\begin{axis}[mcstyle, name=top, height=5.2cm, ymode=log, xmin=0, xmax=250, ymin=5e-3, ymax=1e3, xticklabels={},
  ylabel={$d\sigma/dp_T$ [pb/GeV]}, title={$Z$（$m_{ee}>50$~GeV）}, legend pos=north east]
\addplot[mcblue, thick] table[x=x, y=mc]{figs/mc_pt_z.dat};\addlegendentry{MC@NLO + Pythia 8}
\addplot[ownorange, thick, dashed] table[x=x, y=own]{figs/mc_pt_z.dat};\addlegendentry{LO の $Z$$+$ジェット（自分）}
\end{axis}
\begin{axis}[mcstyle, at={(top.south west)}, anchor=north west, yshift=-0.12cm, height=3.0cm, xmin=0, xmax=250, ymin=0.2, ymax=1.9,
  ytick={0.5,1,1.5}, xlabel={ボソンの $p_T$ [GeV]}, ylabel={比}]
\addplot[gray, thin] coordinates {(0,1) (250,1)};
\addplot[mcblue, only marks, mark=*, mark size=1pt, error bars/.cd, y dir=both, y explicit]
  table[x=x, y=ratio, y error=err]{figs/mc_ptratio_z.dat};
\end{axis}
\end{tikzpicture}\hfill
\begin{tikzpicture}
\begin{axis}[mcstyle, name=top, height=5.2cm, ymode=log, xmin=0, xmax=250, ymin=1e-3, ymax=20, xticklabels={},
  ylabel={$d\sigma/dp_T$ [pb/GeV]}, title={$H$}, legend pos=north east]
\addplot[mcblue, thick] table[x=x, y=mc]{figs/mc_pt_h.dat};\addlegendentry{MC@NLO + Pythia 8}
\addplot[ownorange, thick, dashed] table[x=x, y=own]{figs/mc_pt_h.dat};\addlegendentry{LO の $H$$+$ジェット（自分）}
\end{axis}
\begin{axis}[mcstyle, at={(top.south west)}, anchor=north west, yshift=-0.12cm, height=3.0cm, xmin=0, xmax=250, ymin=0.2, ymax=1.9,
  ytick={0.5,1,1.5}, xlabel={ボソンの $p_T$ [GeV]}, ylabel={比}]
\addplot[gray, thin] coordinates {(0,1) (250,1)};
\addplot[mcblue, only marks, mark=*, mark size=1pt, error bars/.cd, y dir=both, y explicit]
  table[x=x, y=ratio, y error=err]{figs/mc_ptratio_h.dat};
\end{axis}
\end{tikzpicture}
\caption{ボソンの横運動量の分布（5~GeV ごと）。上段は MC@NLO と Pythia 8（実線）と、自分で4次元の実の放射から求めた LO の V$+$ジェット（破線。$H$ は $q\bar q$ を除く）、下段はその比（MC@NLO / 自分の LO）と MC@NLO の統計の誤差。大きい $p_T$ で比は 1 に近づく。LO の V$+$ジェットは $p_T\to0$ で発散するので、最初のビン（0〜5~GeV）には描いていない。}
\label{fig:mcpt}
\end{figure}

\begin{figure}[htbp]
\centering
\begin{tikzpicture}
\begin{axis}[mcstyle, height=5.2cm, xmin=0, xmax=30, ymin=0, ymax=0.12, xlabel={ボソンの $p_T$ [GeV]},
  ylabel={$\frac1\sigma\,d\sigma/dp_T$ [1/GeV]}, title={$Z$（$m_{ee}>50$~GeV）}, legend pos=north east, scaled y ticks=false,
  yticklabel style={/pgf/number format/fixed, /pgf/number format/precision=3}]
\addplot[mcblue, thick] table[x=x, y=mc]{figs/mc_lowpt_z.dat};\addlegendentry{MC@NLO + Pythia 8}
\addplot[ownorange, thick, dashed] table[x=x, y=own]{figs/mc_lowpt_z.dat};\addlegendentry{LO の $Z$$+$ジェット（自分）}
\end{axis}
\end{tikzpicture}\hfill
\begin{tikzpicture}
\begin{axis}[mcstyle, height=5.2cm, xmin=0, xmax=30, ymin=0, ymax=0.05, xlabel={ボソンの $p_T$ [GeV]},
  ylabel={$\frac1\sigma\,d\sigma/dp_T$ [1/GeV]}, title={$H$}, legend pos=north east, scaled y ticks=false,
  yticklabel style={/pgf/number format/fixed, /pgf/number format/precision=3}]
\addplot[mcblue, thick] table[x=x, y=mc]{figs/mc_lowpt_h.dat};\addlegendentry{MC@NLO + Pythia 8}
\addplot[ownorange, thick, dashed] table[x=x, y=own]{figs/mc_lowpt_h.dat};\addlegendentry{LO の $H$$+$ジェット（自分）}
\end{axis}
\end{tikzpicture}
\caption{小さい横運動量（1~GeV ごと、MC@NLO の全断面積で割ったもの）。固定次数の LO の V$+$ジェット（破線）は $p_T\to0$ で発散し、図の上にはみ出す。シャワーをつないだ MC@NLO（実線）は有限で、$H$ は 9〜12~GeV に山を持つ。$Z$ の数 GeV より下は、切ってある primordial $k_T$ とハドロン化に左右される。}
\label{fig:mclow}
\end{figure}

\begin{figure}[htbp]
\centering
\begin{tikzpicture}
\begin{axis}[mcstyle, name=top, height=5.2cm, xmin=-4.5, xmax=4.5, ymin=0, ymax=420, xticklabels={}, ylabel={$d\sigma/dy$ [pb]},
  title={$Z$（$66<m_{ee}<116$~GeV）}, legend pos=north west]
\addplot[mcblue, thick] table[x=x, y=mc]{figs/mc_y_z.dat};\addlegendentry{MC@NLO + Pythia 8}
\addplot[ownorange, thick, dashed] table[x=x, y=ref]{figs/mc_y_z.dat};\addlegendentry{NLO（自分）}
\end{axis}
\begin{axis}[mcstyle, at={(top.south west)}, anchor=north west, yshift=-0.12cm, height=3.0cm, xmin=-4.5, xmax=4.5, ymin=0.8, ymax=1.2,
  ytick={0.9,1,1.1}, xlabel={ボソンのラピディティ $y$}, ylabel={比}]
\addplot[gray, thin] coordinates {(-4.5,1) (4.5,1)};
\addplot[mcblue, only marks, mark=*, mark size=1pt, error bars/.cd, y dir=both, y explicit]
  table[x=x, y=ratio, y error=err]{figs/mc_yratio_z.dat};
\end{axis}
\end{tikzpicture}\hfill
\begin{tikzpicture}
\begin{axis}[mcstyle, name=top, height=5.2cm, xmin=-4.5, xmax=4.5, ymin=0, ymax=9.5, xticklabels={}, ylabel={$d\sigma/dy$ [pb]},
  title={$H$}, legend pos=north west]
\addplot[mcblue, thick] table[x=x, y=mc]{figs/mc_y_h.dat};\addlegendentry{MC@NLO + Pythia 8}
\addplot[foaqua, thick, dotted] table[x=x, y=ref]{figs/mc_y_h.dat};\addlegendentry{固定次数の NLO（MG5）}
\end{axis}
\begin{axis}[mcstyle, at={(top.south west)}, anchor=north west, yshift=-0.12cm, height=3.0cm, xmin=-4.5, xmax=4.5, ymin=0.8, ymax=1.2,
  ytick={0.9,1,1.1}, xlabel={ボソンのラピディティ $y$}, ylabel={比}]
\addplot[gray, thin] coordinates {(-4.5,1) (4.5,1)};
\addplot[mcblue, only marks, mark=*, mark size=1pt, error bars/.cd, y dir=both, y explicit]
  table[x=x, y=ratio, y error=err]{figs/mc_yratio_h.dat};
\end{axis}
\end{tikzpicture}
\caption{ボソンのラピディティの分布。上段は MC@NLO と Pythia 8（実線）と NLO（$Z$ は自分の NLO、破線。$H$ は同じ模型の固定次数の NLO、点線）、下段はその比と誤差（$H$ は両方の統計の誤差）。シャワーはラピディティをほとんど変えない。}
\label{fig:mcy}
\end{figure}

\paragraph{(c) 結果。}表~\ref{tab:mc} と図~\ref{fig:mcpt}〜\ref{fig:mcy} にまとめる。期待 (i)〜(iv) を順に見る。
\begin{description}
\item[(i) 全断面積] MC@NLO の事象の重みの和は、$Z$ の全体で $-0.06\sigma$、$66<m_{ee}<116$~GeV で $-0.33\sigma$、$H$ で $-0.45\sigma$ と、自分の NLO と一致した。シャワーは確率を保つので、MC@NLO の規格化は NLO のものである。
\item[(ii) ラピディティ]（図~\ref{fig:mcy}）$|y|<2.5$ のビンで、MC@NLO と NLO の比は 0.97〜1.02（$Z$、自分の NLO と）、0.97〜1.04（$H$、同じ模型の固定次数の NLO と。$gg\to H$ の NLO のラピディティは自分で導いていない）で、ビンの統計の誤差（1〜2\%）と同じくらいだった。$\chi^2$ は 40 ビンで 60.7、34 ビンで 62.7 で、数 \% の形の違い（$\alpha_s^2$ の違い）が少し見えている。
\item[(iii) 大きい $p_T$]（図~\ref{fig:mcpt}）MC@NLO と自分の LO の V$+$ジェットの比は $p_T$ とともに 1 に近づき、150〜250~GeV で $Z$ が $1.00\pm0.05$、$H$ が $1.035\pm0.018$ になった。途中の 20〜100~GeV では、$Z$ で 1.2〜1.4、$H$ で約 1.5 と大きい。NLO の規格化を持つ Born の事象にシャワーの放射が加わるためと考えられ、$K$ の大きい $H$ ほど差が大きい。
\item[(iv) 小さい $p_T$]（図~\ref{fig:mclow}）固定次数で発散する $p_T\to0$ の分布は有限になった。$H$ は 9〜12~GeV に Sudakov の山を持ち（1~GeV あたり全体の 2.6\%）、5~GeV より下は全体の 6.8\% である。$Z$ は 1〜4~GeV でほぼ平ら（1~GeV あたり全体の 5\%）で、0〜1~GeV に 10\% がある。数 GeV より下は、切ってある primordial $k_T$ やハドロン化に左右される領域なので、比べていない。
\end{description}

\begin{table}[htbp]
\centering\small
\caption{MC@NLO と Pythia 8（パートンの段階、各20万事象）と自分の計算。断面積は pb、括弧は誤差とプル。$p_T$ の比は、MC@NLO を自分の LO の V$+$ジェット、H$+$ジェット（$q\bar q$ を除く）で割ったもの。}
\label{tab:mc}
\begin{tabular}{lll}
\toprule
 & $Z$（$m_{ee}>50$~GeV） & $H$ \\
\midrule
全断面積：自分の NLO & 2064.54 & 27.037 \\
全断面積：MC@NLO & 2064.2(5.5) ($-0.06$) & 27.003(75) ($-0.45$) \\
$66<m_{ee}<116$~GeV：自分、MC@NLO & 1966.68、1964.9(5.4) ($-0.33$) & \\
ラピディティの比（$|y|<2.5$） & 0.973〜1.024（自分の NLO と） & 0.975〜1.039（固定次数の NLO と） \\
$p_T$ の比 5〜20~GeV & 0.813(3) & 0.582(3) \\
\phantom{$p_T$ の比 }20〜50~GeV & 1.358(8) & 1.522(7) \\
\phantom{$p_T$ の比 }50〜100~GeV & 1.182(13) & 1.494(10) \\
\phantom{$p_T$ の比 }100〜150~GeV & 1.040(29) & 1.174(15) \\
\phantom{$p_T$ の比 }150〜250~GeV & 1.001(47) & 1.035(18) \\
$p_T<5$~GeV の割合 & 0.311 & 0.068 \\
\bottomrule
\end{tabular}
\end{table}

\subsection{手順 (4)：SSM の $Z'$ の NLO と質量の下限}\label{sec:step4}

\paragraph{考え方。}$Z'$ は、$Z$ と同じ形でフェルミオンに結合する、重い中性のゲージボソンの仮想的な例である（SSM、sequential standard model）。LHC では Drell--Yan と同じく $q\bar q\to Z'\to\ell^+\ell^-$ で作られ、$m_{\ell\ell}$ の分布に山を作る。計算環境の記録の段階 6~\cite{env} では、自作の模型 \texttt{ZpSSM.fr} で LO の断面積を求め、ATLAS の 139~fb$^{-1}$ のダイレプトンの分布~\cite{atlasll}（参照）から $\sigma\times B(Z'\to\ell\ell)$ の 95\% の上限を出した。SSM の断面積が上限を上回る質量は除外されるので、その境目が質量の下限になる（5.04~TeV）。比べる断面積が大きいほど下限は上がるので、NLO の補正は下限を動かす。$Z'$ はクォークにベクトルと軸性の結合でつくので、手順 (1) の係数関数（質量0のクォークではベクトルと軸性で同じ）がそのまま使え、Born の $\sigma_0$ を $Z'$ のものに替えるだけでよい。手順 (4) は \texttt{Src/run\_zp\_check.sh} で再現できる（約10分。段階 6 の出力が要る）。

\paragraph{(a) 自分の NLO。}\texttt{zp\_hadronic.py} は、
\[
\sigma_0(Q^2)=\frac13\,\frac{Q^2}{48\pi}\,g^4\,\frac{(g_L^{q2}+g_R^{q2})(g_L^{e2}+g_R^{e2})}{(Q^2-M^2)^2+M^2\Gamma^2},\qquad g=\frac e{s_Wc_W},\quad g_L=T_3-Qs_W^2,\quad g_R=-Qs_W^2
\]
（$e^+e^-$ の1つのフレーバー）を式~\eqref{eq:dqq}、\eqref{eq:dqg} の畳み込みに入れる。設定は段階 6 にそろえた。陽子のクォークは $u,d,s,c$（段階 6 では $b$ は質量を持ち陽子に入らない）、NNPDF4.0 NLO、$\mu_R=\mu_F=M$、$\Gamma$ は MadWidth の値（$\Gamma/M\approx0.03$）、$m_{\ell\ell}$ は 10~GeV〜13~TeV である。LO は、段階 6 の自分の積分と9つの質量で $2\times10^{-6}$ 以内、\MGaMC{} の LO と 1.3\% 以内（段階 6 で分かっている裾の積分の差）で一致した。$K=\sigma_{\mathrm{NLO}}/\sigma_{\mathrm{LO}}$ を表~\ref{tab:zp} と図~\ref{fig:zp}（左）に示す。ピークの付近（$|m_{\ell\ell}-M|<2\Gamma$）の $K$ は、0.5~TeV の 1.18 から 6~TeV の 1.41 まで質量とともに大きくなる。$\tau=M^2/S$ が大きいほど、$z\to1$ の軟らかいグルーオンの対数（$\mathcal D_1$）が効くためである。一方、全体の $K$ は重い $Z'$ で下がる。重い $Z'$ では、低い質量の裾（$m\ll M$）が断面積の大きな部分を占め、そこでは $\mu_F=M$ が $m$ から離れて $K$ が小さいからである。9通りのスケールの範囲は $\pm1$〜4\% である。

\paragraph{(b) NLO の UFO と MadLoop。}aMC@NLO で確かめるため、QCD（$u,c,d,s$）、電子、段階 6 と同じ結合の $Z'$ だけの模型 \texttt{QCDZp.fr} を書き、手順 (2) と同じ NLOCT の三つの段（\texttt{zp\_generate\_nlo.sh}、約1分）で NLO の UFO を作った。$Z'$ の頂点は繰り込み可能なので、NLOCT は正しく働く。$q\bar qZ'$ の R2 の頂点は $-2C_F\frac{g_s^2}{16\pi^2}\times$（木の頂点）で、ベクトルと軸性の部分に同じ係数が付いた。紫外の対抗項は0である（質量0のクォークの質量殻の場の繰り込みは0）。\texttt{zp\_madloop.py} の MadLoop（\texttt{u u\textasciitilde{} > e+ e- [virt=QCD]}、$\sqrt{\hat s}=2.5$~TeV、$M=3$~TeV）は、Born が自分の式と $10^{-15}$ で、$A$、$B$、$C$ が手順 (1) の形状因子の $2\,\mathrm{Re}F^{(1)}$（MadLoop の規格化で $A=-4C_F$、$B=-6C_F$、$C=C_F(2\pi^2-16)$）と三つの $\mu_R$ で $5\times10^{-15}$ 以内で一致した。

\paragraph{(c) \MGaMC{} の固定次数。}\texttt{zp\_mg5.sh} は \texttt{p p > e+ e- [QCD]} を 1、3、5~TeV で走らせる（$m_{ee}>225$~GeV、$\mu=M$、9通りのスケールの重み）。LO と NLO の $m_{ee}>225$~GeV の断面積と、ピークの窓の断面積は、9通りのスケールで自分の値と一致した。NLO と LO の比は、1~TeV で 1.2052(13)、1.2063(14)（自分 1.2037、1.2050）、3~TeV で 1.2930(8)、1.3047(6)（1.2925、1.3051）、5~TeV で 1.2978(16)、1.3763(24)（1.2981、1.3739）である（$m_{ee}>225$~GeV、窓の順。図~\ref{fig:zp} の左の点）。初めの 3~TeV の NLO の $m_{ee}>225$~GeV だけは $+2.9\sigma$（$+0.33\%$）ずれた。ずれはピークの外に集まっていたので、種を替え精度を上げてもう一度走らせると $+1.5\sigma$（$+0.034\%$）になった。手順 (1) の $W$ と同じく、\MGaMC{} の表示する誤差は実際のばらつきより小さいことがある。

\paragraph{(d) NLO での質量の下限。}\texttt{zp\_limit.py} は、段階 6 の Delphes の事象（各質量8000事象、ATLAS の card）を、真の $m_{\ell\ell}$ での $K(m)=(\dd\sigma/\dd m)_{\mathrm{NLO}}/(\dd\sigma/\dd m)_{\mathrm{LO}}$ で重み付けし直す。こうして信号の $m_{\ell\ell}$ の型と規格化を NLO にし、段階 6 と同じ尤度と CLs で上限を出し直す（段階 6 の関数をそのまま使う）。真の $m$ には、事象の記録に $Z'$ があればその質量を、なければハードな過程の2つのレプトンの不変質量を使う。\MGaMC{} は、仮想質量がピークから遠い $Z'$ を LHE に書かない。また、$Z'$ が記録にある事象では、崩壊のレプトンは終状態としてしか残らない。事象の重みの平均 $\langle K\rangle$ は、自分の $m_{\ell\ell}$ の全体の $K$ と 1.5\% 以内で合った。重みを 1 にすると段階 6 の下限（観測 5036~GeV、期待 4988~GeV）が再現され、NLO では観測 5204~GeV、期待 5159~GeV になった（表~\ref{tab:zp}、図~\ref{fig:zp} 右）。NLO の補正で下限は約 170~GeV 上がる。ATLAS の値は、NNLO の断面積で 5.1~TeV（参照）である。上限と SSM の断面積の比は、5~TeV で自分の NLO 0.668、ATLAS 0.832 である。この差には、上限そのものの違い（自分の Delphes の効率と型、背景のフィット）と、断面積の次数の違い（ATLAS の NNLO と電弱の補正）が混ざっている。NLO の放射の反跳によるレプトンの受容率の変化は入れていない（事象は LO の運動学のまま）。

\begin{table}[htbp]
\centering\small
\caption{SSM の $Z'$（13~TeV、$e^+e^-$ の1つのフレーバー）。$\sigma_{\mathrm{LO}}$ は自分の LO（$m_{\ell\ell}$ は 10~GeV〜13~TeV）、$K$ は自分の NLO との比（全体、$m_{\ell\ell}>225$~GeV、$|m_{\ell\ell}-M|<2\Gamma$）。右の3列は観測の上限（$ee$ と $\mu\mu$）と SSM の断面積の比で、1 を下回る質量が除外される。ATLAS は NNLO の断面積（参照）。}
\label{tab:zp}
\begin{tabular}{rrrrrrrrr}
\toprule
$M$ [TeV] & $\Gamma/M$ & $\sigma_{\mathrm{LO}}$ [fb] & $K$ & $K_{>225}$ & $K_{\mathrm{peak}}$ & 上限比 LO & 上限比 NLO & ATLAS \\
\midrule
0.5 & 0.0279 & 6816 & 1.174 & 1.177 & 1.177 & 0.000401 & 0.000341 & 0.000222 \\
1.0 & 0.0291 & 464.8 & 1.200 & 1.204 & 1.205 & 0.00127 & 0.00106 & 0.000995 \\
2.0 & 0.0295 & 18.38 & 1.245 & 1.252 & 1.256 & 0.00502 & 0.00398 & 0.00326 \\
3.0 & 0.0296 & 1.699 & 1.273 & 1.293 & 1.305 & 0.0186 & 0.0141 & 0.0139 \\
4.0 & 0.0296 & 0.2502 & 1.264 & 1.309 & 1.345 & 0.176 & 0.129 & 0.173 \\
4.5 & 0.0296 & 0.1154 & 1.244 & 1.306 & 1.360 & 0.358 & 0.259 & 0.318 \\
5.0 & 0.0296 & 0.06031 & 1.221 & 1.298 & 1.374 & 0.932 & 0.668 & 0.832 \\
5.5 & 0.0296 & 0.03508 & 1.199 & 1.289 & 1.389 & 2.52 & 1.79 & 2.53 \\
6.0 & 0.0296 & 0.02221 & 1.182 & 1.282 & 1.408 & 9.56 & 6.78 & 8.23 \\
\midrule
\multicolumn{6}{l}{質量の下限（観測、期待） [GeV]} & 5036、4988 & 5204、5159 & 5100 \\
\bottomrule
\end{tabular}
\end{table}

\begin{figure}[htbp]
\centering
\begin{tikzpicture}[baseline=(current bounding box.north)]
\begin{axis}[mcstyle, height=6.0cm, xmin=0.3, xmax=6.2, ymin=1.1, ymax=1.45, xlabel={$M_{Z'}$ [TeV]}, ylabel={$K=\sigma_{\mathrm{NLO}}/\sigma_{\mathrm{LO}}$},
  title={$K$ 因子}, legend style={at={(0.5,-0.3)}, anchor=north, legend columns=2}, x filter/.expression={x/1000}]
\addplot[mcblue, thick, dotted] table[x=M, y=K]{figs/zp_k.dat};\addlegendentry{自分：全体}
\addplot[mcblue, thick, dashed] table[x=M, y=Kcut]{figs/zp_k.dat};\addlegendentry{自分：$m_{\ell\ell}>225$~GeV}
\addplot[mcblue, thick] table[x=M, y=Kpeak]{figs/zp_k.dat};\addlegendentry{自分：$|m_{\ell\ell}-M|<2\Gamma$}
\addplot[ownorange, only marks, mark=o, mark size=2pt, error bars/.cd, y dir=both, y explicit]
  table[x=M, y=Kcut, y error=Kcuterr]{figs/zp_kmg5.dat};\addlegendentry{\MGaMC{}：$m_{\ell\ell}>225$~GeV}
\addplot[ownorange, only marks, mark=square*, mark size=1.6pt, error bars/.cd, y dir=both, y explicit]
  table[x=M, y=Kpeak, y error=Kpeakerr]{figs/zp_kmg5.dat};\addlegendentry{\MGaMC{}：窓}
\end{axis}
\end{tikzpicture}\hfill
\begin{tikzpicture}[baseline=(current bounding box.north)]
\begin{axis}[mcstyle, height=6.0cm, xmin=0.3, xmax=6.2, ymode=log, ymin=1e-4, ymax=30, xlabel={$M_{Z'}$ [TeV]},
  ylabel={観測の上限 $/\,\sigma_{\mathrm{SSM}}$}, title={上限と SSM の断面積の比}, legend pos=south east, x filter/.expression={x/1000}]
\addplot[gray, thin, forget plot] coordinates {(300,1) (6200,1)};
\addplot[ownorange, thick, dashed, mark=*, mark size=1.2pt] table[x=M, y=lo]{figs/zp_ratio.dat};\addlegendentry{自分（LO の断面積）}
\addplot[mcblue, thick, mark=*, mark size=1.2pt] table[x=M, y=nlo]{figs/zp_ratio.dat};\addlegendentry{自分（NLO の断面積）}
\addplot[foaqua, thick, dotted, mark=triangle*, mark size=1.6pt] table[x=M, y=atlas]{figs/zp_ratio.dat};\addlegendentry{ATLAS（NNLO、参照）}
\end{axis}
\end{tikzpicture}
\caption{SSM の $Z'$。左：$K$ 因子。線は自分の NLO（全体、$m_{\ell\ell}>225$~GeV、ピークの窓）、点は \MGaMC{} の固定次数（1、3、5~TeV、3~TeV は2回目の実行）。右：観測の上限と SSM の断面積の比。1 の線（灰色）を下回る質量が除外され、上と交わる所が質量の下限になる（LO 5.04~TeV、NLO 5.20~TeV、ATLAS 5.1~TeV）。}
\label{fig:zp}
\end{figure}

\section{議論と今後}\label{sec:discussion}

\paragraph{手順 (1) で分かったこと。}第一の問いには肯定で答えられる。$d$ 次元の実の放射と形状因子に $\MSbar$ の共線の対抗項を足すと、$q\bar q$ の2重極と単極、$qg$ の単極が記号として消え、$\ln\mu_R$ が落ち、係数関数は $\MSbar$ のパートン分布の規約（グルーオンの偏極の平均 $1/(2(1-\epsilon))$）と合う形になった。この規約は、$\MSbar$ で決めた NNPDF4.0 と組み合わせたハドロンの断面積が \MGaMC{} と一致したことで、実際に確かめられた。$\delta(1-z)$ の係数 $8\zeta_2-16$ に入る時間的な形状因子の $\pi^2$ の項と、$\mathcal D_1$ の軟らかい放射の対数が $q\bar q$ を約 $+29\%$ 上げ、$qg$ が約 $-15\%$ 下げて、$Z$、$W$ の共鳴の付近で $K\approx1.14$〜$1.15$ になる。

第二の問いにも、手順 (1) の範囲で肯定で答えられる。電弱の結合と Breit--Wigner の伝播関数を持つ \texttt{loop\_sm} の Drell--Yan で、\MGaMC{} の固定次数の LO と NLO は、全断面積、$m_{\ell\ell}$ の分布、ラピディティの分布、9通りの $(\mu_R,\mu_F)$ の重みの付け直しで、自分の畳み込みと統計の誤差の範囲で一致した（$W$ の NLO の全体は、$m_{e\nu}<1$~GeV の打ち消しによる積分のばらつきの範囲で）。QCD ノートの $t\bar t$ と合わせて、\texttt{[QCD]} の FKS の引き算は、始状態の $q\bar q$、$qg$、$gg$ のチャンネルで確かめたことになる。比べる途中で、自分の積分の Breit--Wigner の変数だけの求積が $W$ の裾を $10^{-3}$ 少なく見積もっていたことが分かった。点を倍にする確かめは、比べる範囲そのものでする必要がある（計算環境の記録~\cite{env} の表に加えた）。

\paragraph{手順 (2) で分かったこと。}$gg\to H$ の係数関数は、Drell--Yan と同じ手順で $d$ 次元で求まり、三つのチャンネルが文献と記号として一致した。\MGaMC{} の側では、ツールの使い方に三つの注意が要った。次元5の演算子の NLO の UFO は NLOCT だけでは作れないこと（R2 と紫外の対抗項を自分で導いて入れると作れ、MadLoop は直しなしで正しく、aMC@NLO は次数の明示と頂点のないダミーの粒子を足すだけで動いた）、配布の HC\_NLO\_X0 は MG5 3.8.1 で軽いクォークの対抗項の極が落ちること（極の確かめで止まる）、その結合の次数のために $\mu_R$ の重みの付け直しが効かないことである。さらに、aMC@NLO の \texttt{[QCD]} は FKS の対に写らない有限な実の放射（$q\bar q\to Hg$、LO の $0.13\%$）を含まない。これらは自分の計算と並べて初めて分かったことで、計算環境の記録~\cite{env} の表に加えた。物理としては、$K=2.29$ と大きく、スケールの範囲が NLO でも縮まないので、$gg\to H$ の精密な予言には NNLO 以上が要る。

\paragraph{手順 (3) で分かったこと。}第三の問いには次のように答えられる。MC@NLO で Pythia 8 のシャワーをつないでも、全断面積は NLO のまま変わらず（$Z$、$H$ とも統計の範囲）、ボソンのラピディティは数 \% しか変わらない。変わるのは $p_T$ の分布で、固定次数の NLO が $p_T\to0$ で発散するところが有限な分布になり、大きい $p_T$ では LO の V$+$ジェットに戻る。途中の 20〜100~GeV では、MC@NLO は LO の V$+$ジェットより $Z$ で 20〜40\%、$H$ で約 50\% 大きく、この領域の予言は NLO の補正の大きさとシャワーのつなぎ方に敏感である。ツールの側では、スケールの重みを切った MC@NLO のシャワーがメモリを使い尽くすこと、Pythia 8 のリンクに gfortran の場所が要ること、QED の次数を持たない模型で事象の生成が止まることが分かり、計算環境の記録~\cite{env} の表に加えた。

\paragraph{手順 (4) で分かったこと。}Drell--Yan の係数関数はベクトルと軸性の流れで同じなので、新しいボソンの NLO は Born を替えるだけで求まる。$Z'$ のピークの $K$ は質量とともに大きくなり（6~TeV で 1.41）、全体の $K$（低い質量の裾を含む）とは 20\% 近く違う。上限の計算では、信号の $m_{\ell\ell}$ の型と規格化の両方に $K(m)$ を入れる必要があった。下限は約 170~GeV 上がり、ATLAS の NNLO の値（参照）を上回った。この差は、断面積の次数だけでなく、段階 6 の上限そのもの（Delphes の効率と型）にもよる。NLOCT は、繰り込み可能な $Z'$ の頂点では正しい R2 を出し、MadLoop と aMC@NLO は直しなしで動いた。

\paragraph{今後。}手順 (3) では、ハドロン化と多重パートン相互作用、primordial $k_T$ を切ったパートンの段階で比べた。数 GeV より下の $p_T$ の形と、ジェットの観測量は比べていない。手順 (4) の質量の下限には、NLO の放射の反跳によるレプトンの受容率の変化と、NNLO と電弱の補正を入れていない。ラピディティの分布を除く、レプトンにカットを掛けた分布（fiducial な断面積）は、自分の計算がレプトンの角度を積分しているので比べていない。$gg\to H$ の NLO のラピディティの分布と、トップの質量を残した NLO（2ループの仮想の補正）も比べていない。

\appendix
\section{Drell--Yan の $d$ 次元の導出の細部}\label{app:dy}

\paragraph{位相空間。}$q\bar q\to V^*(Q^2)+$質量0のパートンの $d$ 次元の2体の位相空間は、パートンの運動量 $|\boldsymbol k|=\frac{\sqrt{\hat s}}2(1-z)$、$p_1$ との角 $\theta$、$v=\frac{1-\cos\theta}2$ で
\[
\dd\Phi_2=\frac{|\boldsymbol k|^{d-3}}{4(2\pi)^{d-2}\sqrt{\hat s}}\,\dd\Omega_{d-2},\qquad
\dd\Omega_{d-2}=\frac{2\pi^{1-\epsilon}}{\Gamma(1-\epsilon)}\,2\cdot4^{-\epsilon}[v(1-v)]^{-\epsilon}\dd v
\]
で、まとめると第~\ref{sec:coef} 節の式になる。Born の1体の位相空間は $2\pi\delta(\hat s-Q^2)=\frac{2\pi}{\hat s}\delta(1-z)$ で、レプトンの部分は Born と実の放射で共通である。

\paragraph{$v$ の積分。}式~\eqref{eq:real} を $v$ について部分分数にすると、$v^n$ と $(1-v)^m$ の和になり、$\int_0^1\dd v\,v^{n-\epsilon}(1-v)^{-\epsilon}=\frac{\Gamma(n+1-\epsilon)\Gamma(1-\epsilon)}{\Gamma(n+2-2\epsilon)}$ で積分できる。$q\bar q$ の $(1-z)^2\times$（積分）は $z=1$ で正則で、$z=1$ での値は $8C_Fe^{\epsilon(\gamma_E+\ln(\mu_R^2/Q^2))}\Gamma(-\epsilon)/\Gamma(1-2\epsilon)$ である。

\paragraph{分布の展開。}$H(z,\epsilon)=H_{-1}/\epsilon+H_0+H_1\epsilon$ を正則な部分として $H(z,\epsilon)(1-z)^{-1-2\epsilon}=-\frac{H(1,\epsilon)}{2\epsilon}\delta(1-z)+H(z,\epsilon)\{[\frac1{1-z}]_+-2\epsilon[\frac{\ln(1-z)}{1-z}]_+\}+O(\epsilon)$ とし、$c(z)[g]_+=c(1)[g]_++(c(z)-c(1))g(z)$ で正則な部分に分ける。$\delta(1-z)$ には $H(1,\epsilon)$ の $O(\epsilon)$ まで、$[\frac{\ln(1-z)}{1-z}]_+$ には $H_{-1}$ だけが要る。

\section{ラピディティの分布の NLO}\label{app:rap}

$v$ を固定すると、レプトンの対のパートンの重心系でのラピディティは $y^*=\frac12\ln\frac{z+v(1-z)}{1-v(1-z)}$ で、$v=0$（$k\parallel p_1$）で $\frac12\ln z$、$v=1$ で $-\frac12\ln z$、$z=1$ で0である。ハドロンの系では $y=\ln x_1+\frac12\ln(z/\tau)+y^*$ なので、観測量（ラピディティのビン）で重みを付けた光度
\[
G(z,v)=\int\dd t\,x f_i(e^t)\,x f_j(\tau e^{-t}/z)\,\Theta\bigl(t+\tfrac12\ln(z/\tau)+y^*(z,v)\in\text{ビン}\bigr)
\]
は $v$ の連続関数で、$G_a=G(z,0)$、$G_b=G(z,1)$、$G(1)=G(1,v)$ と書く。$q\bar q$ では $R(z,v,\epsilon)=(1-z)^2v(1-v)r_{q\bar q}=R_0+\epsilon R_1$、$[v(1-v)]^{-1-\epsilon}=v^{-1-\epsilon}(1-v)^{-\epsilon}+v^{-\epsilon}(1-v)^{-1-\epsilon}$ として、$z<1$ で $v^{-1-\epsilon}=-\delta(v)/\epsilon+[1/v]_++\cdots$ と展開する。$v=0$、1 の極は同じ運動学の共線の対抗項と、$\delta(1-z)$ の部分は式~\eqref{eq:dqq} の $\delta$ の係数 $c_\delta=C_F(8\zeta_2-16)-6C_FL_F$ と同じものになり
\begin{align*}
[\Delta_{q\bar q}G]^{(1)}&=c_\delta G(1)+\int_\tau^1\frac{\dd z}{1-z}\Bigl\{2C_F(1-z)^2(G_a+G_b)-2C_F(1+z^2)\ln z\,(G_a+G_b)\\
&\qquad+\int_0^1\dd v\Bigl[\frac{R_0G-R_0(z,0)G_a}v+\frac{R_0G-R_0(z,1)G_b}{1-v}\Bigr]\Bigr\}\\
&\quad+4C_F\,(1+z^2)\mathcal D_1\otimes(G_a+G_b)-2C_FL_F\,(1+z^2)\mathcal D_0\otimes(G_a+G_b)
\end{align*}
になる（$R_0=2C_F[(1-z)^2(v^2+(1-v)^2)+2z]$、$R_1(z,0)=R_1(z,1)=-2C_F(1-z)^2$、$\mathcal D_k=[\ln^k(1-z)/(1-z)]_+$）。$q(p_1)g(p_2)$ では $S=(1-z)(1-v)r_{qg}=S_0+\epsilon S_1$ がクォークがグルーオンに沿う $v=1$ でだけ特異で
\[
[\Delta_{qg}G]^{(1)}=\int_\tau^1\dd z\Bigl\{2T_RP(z)\bigl(2\ln(1-z)-\ln z\bigr)G_b+4T_Rz(1-z)G_b+\int_0^1\dd v\,\frac{S_0G-S_0(z,1)G_b}{1-v}-2T_RL_FP(z)G_b\Bigr\}
\]
（$P=z^2+(1-z)^2$、$S_0=2T_R[(1-z)^2(1-v)^2+1-2z(1-z)v]$）で、$g(p_1)q(p_2)$ は $v\to1-v$、$G_b\to G_a$ である。$G$ が $v$ によらないとき（$y$ で積分したとき）、両方が式~\eqref{eq:dqq}、\eqref{eq:dqg} に記号として戻ることを \texttt{dy\_feyncalc.wl} の (d2) で確かめた。

\begin{thebibliography}{99}
\bibitem{smnote} 標準模型：$\MSbar$ の電弱の1ループの繰り込みと検算（$V_{\mathrm{CKM}}=1$、FJ のタドポール、一般の $R_\xi$ ゲージ、$d=4$）（本リポジトリ \nolinkurl{models/sm/ResearchNote/note.pdf}）。
\bibitem{ckmnote} 標準模型と CKM 行列：$\MSbar$ の繰り込みと世代を変える過程（FJ のタドポール、一般の $R_\xi$ ゲージ、$d=4$）（本リポジトリ \nolinkurl{models/sm_ckm/ResearchNote/note.pdf}）。
\bibitem{nlonote} 標準模型の NLO の物理量：実の放射を含む崩壊幅と電弱の精密測定（本リポジトリ \nolinkurl{models/sm_nlo/ResearchNote/note.pdf}）。
\bibitem{higgsnote} 標準模型のヒッグスの崩壊：NLO の部分幅と分岐比（本リポジトリ \nolinkurl{models/sm_higgs/ResearchNote/note.pdf}）。
\bibitem{qcdnote} QCD：1ループ・2ループの構造、インスタントンと検算（本リポジトリ \nolinkurl{models/qcd/ResearchNote/note.pdf}）。
\bibitem{aem79} G. Altarelli, R. K. Ellis, G. Martinelli, ``Large perturbative corrections to the Drell-Yan process in QCD,'' Nucl. Phys. B \textbf{157}, 461 (1979)。$\MSbar$ の形は R. K. Ellis, W. J. Stirling, B. R. Webber, \textit{QCD and Collider Physics}（Cambridge University Press, 1996）の 9.2 節による（参照）。
\bibitem{dawson91} S. Dawson, ``Radiative corrections to Higgs boson production,'' Nucl. Phys. B \textbf{359}, 283 (1991)（参照）。
\bibitem{dsz91} A. Djouadi, M. Spira, P. M. Zerwas, ``Production of Higgs bosons in proton colliders: QCD corrections,'' Phys. Lett. B \textbf{264}, 440 (1991)（参照）。
\bibitem{hk02} R. V. Harlander, W. B. Kilgore, ``Next-to-next-to-leading order Higgs production at hadron colliders,'' Phys. Rev. Lett. \textbf{88}, 201801 (2002), arXiv:hep-ph/0201206（参照）。
\bibitem{hcnlo} P. Artoisenet \textit{et al.}, ``A framework for Higgs characterisation,'' JHEP \textbf{11}, 043 (2013), arXiv:1306.6464；F. Maltoni, K. Mawatari, M. Zaro, Eur. Phys. J. C \textbf{74}, 2710 (2014), arXiv:1311.1829；F. Demartin, F. Maltoni, K. Mawatari, B. Page, M. Zaro, Eur. Phys. J. C \textbf{74}, 3065 (2014), arXiv:1407.5089（HC\_NLO\_X0 の版 1.3 の README が挙げる文献）。
\bibitem{yr4ggh} D. de Florian \textit{et al.} (LHC Higgs Cross Section Working Group), arXiv:1610.07922、$gg\to H$ の13~TeV の表（参照）。
\bibitem{nnpdf40} R. D. Ball \textit{et al.} (NNPDF Collaboration), ``The path to proton structure at 1\% accuracy,'' Eur. Phys. J. C \textbf{82}, 428 (2022), arXiv:2109.02653。
\bibitem{fks} S. Frixione, Z. Kunszt, A. Signer, ``Three-jet cross sections to next-to-leading order,'' Nucl. Phys. B \textbf{467}, 399 (1996), arXiv:hep-ph/9512328；R. Frederix, S. Frixione, F. Maltoni, T. Stelzer, ``Automation of next-to-leading order computations in QCD: the FKS subtraction,'' JHEP \textbf{10}, 003 (2009), arXiv:0908.4272。
\bibitem{mg5} J. Alwall \textit{et al.}, ``The automated computation of tree-level and next-to-leading order differential cross sections, and their matching to parton shower simulations,'' JHEP \textbf{07}, 079 (2014), arXiv:1405.0301。
\bibitem{mcatnlo} S. Frixione, B. R. Webber, ``Matching NLO QCD computations and parton shower simulations,'' JHEP \textbf{06}, 029 (2002), arXiv:hep-ph/0204244.
\bibitem{pythia83} C. Bierlich \textit{et al.}, ``A comprehensive guide to the physics and usage of PYTHIA 8.3,'' SciPost Phys.\ Codebases \textbf{8} (2022), arXiv:2203.11601.
\bibitem{atlasll} ATLAS Collaboration, ``Search for high-mass dilepton resonances using 139 fb$^{-1}$ of $pp$ collision data collected at $\sqrt s=13$ TeV with the ATLAS detector,'' Phys.\ Lett.\ B \textbf{796}, 68 (2019), arXiv:1903.06248；HEPData ins1725190（参照）。
\bibitem{atlaswz} ATLAS Collaboration, ``Measurement of $W^\pm$ and $Z$-boson production cross sections in $pp$ collisions at $\sqrt s=13$~TeV with the ATLAS detector,'' Phys. Lett. B \textbf{759}, 601 (2016), arXiv:1603.09222（参照）。
\bibitem{env} ループ計算環境の構築記録（本リポジトリ \nolinkurl{environment/environment.pdf}）。
\end{thebibliography}

\end{document}
